% documentation_FR.tex — WAC (Windows Artefacts Collector) user manual, French edition.
% The English edition is documentation_EN.tex, in this folder: both are maintained together,
% section for section.
%
% Build (XeLaTeX, two passes for the cover page and the table of contents):
%     latexmk -xelatex documentation_FR.tex
% The cover page is ../cover_FR.png.

\documentclass[11pt,a4paper,oneside]{report}

% ---------------------------------------------------------------- language and labels
\usepackage[french]{babel}
\newcommand{\wacdoctitle}{Documentation utilisateur}
\newcommand{\wacpageof}{sur}
\newcommand{\wacnotetitle}{À retenir}
\newcommand{\wacwarningtitle}{Attention}
% Layout shared with every WAC document.
% wac-style.tex — layout shared by every WAC document (user manuals, testing
% guides), French and English editions alike: one copy, so that they cannot
% drift apart.
%
% A document sets, BEFORE % wac-style.tex — layout shared by every WAC document (user manuals, testing
% guides), French and English editions alike: one copy, so that they cannot
% drift apart.
%
% A document sets, BEFORE % wac-style.tex — layout shared by every WAC document (user manuals, testing
% guides), French and English editions alike: one copy, so that they cannot
% drift apart.
%
% A document sets, BEFORE \input{../wac-style.tex}:
%   \usepackage[<language>]{babel}
%   \newcommand{\wacdoctitle}{...}      running header: "WAC — <title>"
%   \newcommand{\wacpageof}{...}        footer: "Page N <of> M"
%   \newcommand{\wacnotetitle}{...}     default title of the "note" box
%   \newcommand{\wacwarningtitle}{...}  default title of the "attention" box
% and AFTER it its \hypersetup (PDF title, author, subject).

% ---------------------------------------------------------------- fonts, language
\usepackage{fontspec}
\setmainfont{Noto Sans}
\setmonofont{DejaVu Sans Mono}[Scale=0.84]

% ---------------------------------------------------------------- layout
\usepackage[top=26mm,bottom=26mm,left=24mm,right=22mm,headheight=14pt]{geometry}
\usepackage{graphicx}
\graphicspath{{../}}
\usepackage{xcolor}
\usepackage{tikz}
\usepackage{booktabs,tabularx,longtable,array}
\usepackage{enumitem}
\usepackage{fancyhdr}
\usepackage{titlesec}
\usepackage[most]{tcolorbox}
\usepackage{listings}
\usepackage{url}
\usepackage{lastpage}
\usepackage[hidelinks,unicode,pdfpagelabels]{hyperref}

% Colours sampled from the cover page
\definecolor{marine}{HTML}{0A1A33}
\definecolor{bleu}{HTML}{2F6DB5}
\definecolor{accent}{HTML}{1F66D1}
\definecolor{fond}{HTML}{EEF3FA}
\definecolor{alerte}{HTML}{B3261E}
\definecolor{fondalerte}{HTML}{FBEDEC}
\definecolor{gris}{HTML}{5B6573}

\setlength{\parindent}{0pt}
\setlength{\parskip}{0.55em}
\linespread{1.08}
\setlist{itemsep=0.25em,topsep=0.35em}
\urlstyle{tt}
\renewcommand{\arraystretch}{1.25}
% Left-aligned table columns: justified narrow columns opened wide gaps
% between words.
\newcolumntype{L}[1]{>{\raggedright\arraybackslash}p{#1}}
% Monospaced file and field names cannot be hyphenated: let the line's spaces
% stretch instead.
\setlength{\emergencystretch}{3em}

% Headings
\titleformat{\chapter}[display]
  {\color{marine}\bfseries\raggedright}
  {\color{bleu}\large\MakeUppercase{\chaptertitlename}~\thechapter}
  {0.3em}
  {\Huge}[{\color{accent}\titlerule[2pt]}]
\titlespacing*{\chapter}{0pt}{0pt}{1.6em}
\titleformat{\section}{\color{marine}\Large\bfseries}{\color{bleu}\thesection}{0.8em}{}
\titleformat{\subsection}{\color{marine}\large\bfseries}{\color{bleu}\thesubsection}{0.7em}{}
\titleformat{\subsubsection}{\color{bleu}\normalsize\bfseries}{}{0em}{}

% Headers and footers
\pagestyle{fancy}
\fancyhf{}
\fancyhead[L]{\small\color{gris}WAC — \wacdoctitle}
\fancyhead[R]{\small\color{gris}\leftmark}
% Footer: "Page N of M" (the word set by \wacpageof), readable. A bare, small, light-grey number went
% unnoticed.
\newcommand{\piedpage}{\small\color{marine}\textbf{Page \thepage}\ \wacpageof\ \pageref*{LastPage}}
\fancyfoot[C]{\piedpage}
\renewcommand{\footrulewidth}{0.4pt}
\renewcommand{\footrule}{\hbox to\headwidth{\color{fond}\leaders\hrule height \footrulewidth\hfill}}
\renewcommand{\headrulewidth}{0.4pt}
\renewcommand{\headrule}{\hbox to\headwidth{\color{fond}\leaders\hrule height \headrulewidth\hfill}}
\renewcommand{\chaptermark}[1]{\markboth{#1}{}}
\fancypagestyle{plain}{\fancyhf{}\fancyfoot[C]{\piedpage}\renewcommand{\headrulewidth}{0pt}}

% Boxes
\newtcolorbox{note}[1][\wacnotetitle]{enhanced,breakable,colback=fond,colframe=bleu,
  boxrule=0pt,leftrule=3pt,arc=0pt,left=8pt,right=8pt,top=6pt,bottom=6pt,
  title={\bfseries #1},coltitle=bleu,colbacktitle=fond,attach title to upper={\par\smallskip}}
\newtcolorbox{attention}[1][\wacwarningtitle]{enhanced,breakable,colback=fondalerte,colframe=alerte,
  boxrule=0pt,leftrule=3pt,arc=0pt,left=8pt,right=8pt,top=6pt,bottom=6pt,
  title={\bfseries #1},coltitle=alerte,colbacktitle=fondalerte,attach title to upper={\par\smallskip}}

% Code
\lstset{basicstyle=\ttfamily\small,columns=fullflexible,keepspaces=true,
  backgroundcolor=\color{fond},frame=l,framerule=3pt,rulecolor=\color{bleu},
  xleftmargin=9pt,framexleftmargin=6pt,aboveskip=0.8em,belowskip=0.8em,
  breaklines=true,breakatwhitespace=false,showstringspaces=false,
  literate={é}{{\'e}}1 {è}{{\`e}}1 {à}{{\`a}}1 {ê}{{\^e}}1 {ô}{{\^o}}1 {ç}{{\c c}}1
           {É}{{\'E}}1 {—}{{---}}1 {’}{{'}}1 {…}{{...}}1 {«}{{\guillemotleft}}1 {»}{{\guillemotright}}1}

% Shortcuts
\newcommand{\wac}{\textsc{wac}}
\newcommand{\opt}[1]{\texttt{-{}-#1}}
\newcommand{\champ}[1]{\texttt{#1}}
\newcommand{\fichier}[1]{\texttt{#1}}

% ============================================================ COVER PAGE
% The cover images are 2:3 (1024 x 1536), taller than an A4 page once widened
% to it: centred, they lost 9 mm at the top — the CEC logo — and 9 mm at the
% bottom. Raised 4 mm above the top of the page instead: the logo, which starts
% 6 mm below the image's edge, stays whole, and the bottom row of text keeps
% 7 mm to the page's edge; only a decorative corner below it is cut.
% Usage: \coverpage{cover_FR.png}
\newcommand{\coverpage}[1]{%
  \begin{titlepage}
  \thispagestyle{empty}
  \begin{tikzpicture}[remember picture,overlay]
    \node[inner sep=0pt,anchor=north,yshift=4mm] at (current page.north)
      {\includegraphics[width=\paperwidth]{#1}};
  \end{tikzpicture}
  \end{titlepage}}
:
%   \usepackage[<language>]{babel}
%   \newcommand{\wacdoctitle}{...}      running header: "WAC — <title>"
%   \newcommand{\wacpageof}{...}        footer: "Page N <of> M"
%   \newcommand{\wacnotetitle}{...}     default title of the "note" box
%   \newcommand{\wacwarningtitle}{...}  default title of the "attention" box
% and AFTER it its \hypersetup (PDF title, author, subject).

% ---------------------------------------------------------------- fonts, language
\usepackage{fontspec}
\setmainfont{Noto Sans}
\setmonofont{DejaVu Sans Mono}[Scale=0.84]

% ---------------------------------------------------------------- layout
\usepackage[top=26mm,bottom=26mm,left=24mm,right=22mm,headheight=14pt]{geometry}
\usepackage{graphicx}
\graphicspath{{../}}
\usepackage{xcolor}
\usepackage{tikz}
\usepackage{booktabs,tabularx,longtable,array}
\usepackage{enumitem}
\usepackage{fancyhdr}
\usepackage{titlesec}
\usepackage[most]{tcolorbox}
\usepackage{listings}
\usepackage{url}
\usepackage{lastpage}
\usepackage[hidelinks,unicode,pdfpagelabels]{hyperref}

% Colours sampled from the cover page
\definecolor{marine}{HTML}{0A1A33}
\definecolor{bleu}{HTML}{2F6DB5}
\definecolor{accent}{HTML}{1F66D1}
\definecolor{fond}{HTML}{EEF3FA}
\definecolor{alerte}{HTML}{B3261E}
\definecolor{fondalerte}{HTML}{FBEDEC}
\definecolor{gris}{HTML}{5B6573}

\setlength{\parindent}{0pt}
\setlength{\parskip}{0.55em}
\linespread{1.08}
\setlist{itemsep=0.25em,topsep=0.35em}
\urlstyle{tt}
\renewcommand{\arraystretch}{1.25}
% Left-aligned table columns: justified narrow columns opened wide gaps
% between words.
\newcolumntype{L}[1]{>{\raggedright\arraybackslash}p{#1}}
% Monospaced file and field names cannot be hyphenated: let the line's spaces
% stretch instead.
\setlength{\emergencystretch}{3em}

% Headings
\titleformat{\chapter}[display]
  {\color{marine}\bfseries\raggedright}
  {\color{bleu}\large\MakeUppercase{\chaptertitlename}~\thechapter}
  {0.3em}
  {\Huge}[{\color{accent}\titlerule[2pt]}]
\titlespacing*{\chapter}{0pt}{0pt}{1.6em}
\titleformat{\section}{\color{marine}\Large\bfseries}{\color{bleu}\thesection}{0.8em}{}
\titleformat{\subsection}{\color{marine}\large\bfseries}{\color{bleu}\thesubsection}{0.7em}{}
\titleformat{\subsubsection}{\color{bleu}\normalsize\bfseries}{}{0em}{}

% Headers and footers
\pagestyle{fancy}
\fancyhf{}
\fancyhead[L]{\small\color{gris}WAC — \wacdoctitle}
\fancyhead[R]{\small\color{gris}\leftmark}
% Footer: "Page N of M" (the word set by \wacpageof), readable. A bare, small, light-grey number went
% unnoticed.
\newcommand{\piedpage}{\small\color{marine}\textbf{Page \thepage}\ \wacpageof\ \pageref*{LastPage}}
\fancyfoot[C]{\piedpage}
\renewcommand{\footrulewidth}{0.4pt}
\renewcommand{\footrule}{\hbox to\headwidth{\color{fond}\leaders\hrule height \footrulewidth\hfill}}
\renewcommand{\headrulewidth}{0.4pt}
\renewcommand{\headrule}{\hbox to\headwidth{\color{fond}\leaders\hrule height \headrulewidth\hfill}}
\renewcommand{\chaptermark}[1]{\markboth{#1}{}}
\fancypagestyle{plain}{\fancyhf{}\fancyfoot[C]{\piedpage}\renewcommand{\headrulewidth}{0pt}}

% Boxes
\newtcolorbox{note}[1][\wacnotetitle]{enhanced,breakable,colback=fond,colframe=bleu,
  boxrule=0pt,leftrule=3pt,arc=0pt,left=8pt,right=8pt,top=6pt,bottom=6pt,
  title={\bfseries #1},coltitle=bleu,colbacktitle=fond,attach title to upper={\par\smallskip}}
\newtcolorbox{attention}[1][\wacwarningtitle]{enhanced,breakable,colback=fondalerte,colframe=alerte,
  boxrule=0pt,leftrule=3pt,arc=0pt,left=8pt,right=8pt,top=6pt,bottom=6pt,
  title={\bfseries #1},coltitle=alerte,colbacktitle=fondalerte,attach title to upper={\par\smallskip}}

% Code
\lstset{basicstyle=\ttfamily\small,columns=fullflexible,keepspaces=true,
  backgroundcolor=\color{fond},frame=l,framerule=3pt,rulecolor=\color{bleu},
  xleftmargin=9pt,framexleftmargin=6pt,aboveskip=0.8em,belowskip=0.8em,
  breaklines=true,breakatwhitespace=false,showstringspaces=false,
  literate={é}{{\'e}}1 {è}{{\`e}}1 {à}{{\`a}}1 {ê}{{\^e}}1 {ô}{{\^o}}1 {ç}{{\c c}}1
           {É}{{\'E}}1 {—}{{---}}1 {’}{{'}}1 {…}{{...}}1 {«}{{\guillemotleft}}1 {»}{{\guillemotright}}1}

% Shortcuts
\newcommand{\wac}{\textsc{wac}}
\newcommand{\opt}[1]{\texttt{-{}-#1}}
\newcommand{\champ}[1]{\texttt{#1}}
\newcommand{\fichier}[1]{\texttt{#1}}

% ============================================================ COVER PAGE
% The cover images are 2:3 (1024 x 1536), taller than an A4 page once widened
% to it: centred, they lost 9 mm at the top — the CEC logo — and 9 mm at the
% bottom. Raised 4 mm above the top of the page instead: the logo, which starts
% 6 mm below the image's edge, stays whole, and the bottom row of text keeps
% 7 mm to the page's edge; only a decorative corner below it is cut.
% Usage: \coverpage{cover_FR.png}
\newcommand{\coverpage}[1]{%
  \begin{titlepage}
  \thispagestyle{empty}
  \begin{tikzpicture}[remember picture,overlay]
    \node[inner sep=0pt,anchor=north,yshift=4mm] at (current page.north)
      {\includegraphics[width=\paperwidth]{#1}};
  \end{tikzpicture}
  \end{titlepage}}
:
%   \usepackage[<language>]{babel}
%   \newcommand{\wacdoctitle}{...}      running header: "WAC — <title>"
%   \newcommand{\wacpageof}{...}        footer: "Page N <of> M"
%   \newcommand{\wacnotetitle}{...}     default title of the "note" box
%   \newcommand{\wacwarningtitle}{...}  default title of the "attention" box
% and AFTER it its \hypersetup (PDF title, author, subject).

% ---------------------------------------------------------------- fonts, language
\usepackage{fontspec}
\setmainfont{Noto Sans}
\setmonofont{DejaVu Sans Mono}[Scale=0.84]

% ---------------------------------------------------------------- layout
\usepackage[top=26mm,bottom=26mm,left=24mm,right=22mm,headheight=14pt]{geometry}
\usepackage{graphicx}
\graphicspath{{../}}
\usepackage{xcolor}
\usepackage{tikz}
\usepackage{booktabs,tabularx,longtable,array}
\usepackage{enumitem}
\usepackage{fancyhdr}
\usepackage{titlesec}
\usepackage[most]{tcolorbox}
\usepackage{listings}
\usepackage{url}
\usepackage{lastpage}
\usepackage[hidelinks,unicode,pdfpagelabels]{hyperref}

% Colours sampled from the cover page
\definecolor{marine}{HTML}{0A1A33}
\definecolor{bleu}{HTML}{2F6DB5}
\definecolor{accent}{HTML}{1F66D1}
\definecolor{fond}{HTML}{EEF3FA}
\definecolor{alerte}{HTML}{B3261E}
\definecolor{fondalerte}{HTML}{FBEDEC}
\definecolor{gris}{HTML}{5B6573}

\setlength{\parindent}{0pt}
\setlength{\parskip}{0.55em}
\linespread{1.08}
\setlist{itemsep=0.25em,topsep=0.35em}
\urlstyle{tt}
\renewcommand{\arraystretch}{1.25}
% Left-aligned table columns: justified narrow columns opened wide gaps
% between words.
\newcolumntype{L}[1]{>{\raggedright\arraybackslash}p{#1}}
% Monospaced file and field names cannot be hyphenated: let the line's spaces
% stretch instead.
\setlength{\emergencystretch}{3em}

% Headings
\titleformat{\chapter}[display]
  {\color{marine}\bfseries\raggedright}
  {\color{bleu}\large\MakeUppercase{\chaptertitlename}~\thechapter}
  {0.3em}
  {\Huge}[{\color{accent}\titlerule[2pt]}]
\titlespacing*{\chapter}{0pt}{0pt}{1.6em}
\titleformat{\section}{\color{marine}\Large\bfseries}{\color{bleu}\thesection}{0.8em}{}
\titleformat{\subsection}{\color{marine}\large\bfseries}{\color{bleu}\thesubsection}{0.7em}{}
\titleformat{\subsubsection}{\color{bleu}\normalsize\bfseries}{}{0em}{}

% Headers and footers
\pagestyle{fancy}
\fancyhf{}
\fancyhead[L]{\small\color{gris}WAC — \wacdoctitle}
\fancyhead[R]{\small\color{gris}\leftmark}
% Footer: "Page N of M" (the word set by \wacpageof), readable. A bare, small, light-grey number went
% unnoticed.
\newcommand{\piedpage}{\small\color{marine}\textbf{Page \thepage}\ \wacpageof\ \pageref*{LastPage}}
\fancyfoot[C]{\piedpage}
\renewcommand{\footrulewidth}{0.4pt}
\renewcommand{\footrule}{\hbox to\headwidth{\color{fond}\leaders\hrule height \footrulewidth\hfill}}
\renewcommand{\headrulewidth}{0.4pt}
\renewcommand{\headrule}{\hbox to\headwidth{\color{fond}\leaders\hrule height \headrulewidth\hfill}}
\renewcommand{\chaptermark}[1]{\markboth{#1}{}}
\fancypagestyle{plain}{\fancyhf{}\fancyfoot[C]{\piedpage}\renewcommand{\headrulewidth}{0pt}}

% Boxes
\newtcolorbox{note}[1][\wacnotetitle]{enhanced,breakable,colback=fond,colframe=bleu,
  boxrule=0pt,leftrule=3pt,arc=0pt,left=8pt,right=8pt,top=6pt,bottom=6pt,
  title={\bfseries #1},coltitle=bleu,colbacktitle=fond,attach title to upper={\par\smallskip}}
\newtcolorbox{attention}[1][\wacwarningtitle]{enhanced,breakable,colback=fondalerte,colframe=alerte,
  boxrule=0pt,leftrule=3pt,arc=0pt,left=8pt,right=8pt,top=6pt,bottom=6pt,
  title={\bfseries #1},coltitle=alerte,colbacktitle=fondalerte,attach title to upper={\par\smallskip}}

% Code
\lstset{basicstyle=\ttfamily\small,columns=fullflexible,keepspaces=true,
  backgroundcolor=\color{fond},frame=l,framerule=3pt,rulecolor=\color{bleu},
  xleftmargin=9pt,framexleftmargin=6pt,aboveskip=0.8em,belowskip=0.8em,
  breaklines=true,breakatwhitespace=false,showstringspaces=false,
  literate={é}{{\'e}}1 {è}{{\`e}}1 {à}{{\`a}}1 {ê}{{\^e}}1 {ô}{{\^o}}1 {ç}{{\c c}}1
           {É}{{\'E}}1 {—}{{---}}1 {’}{{'}}1 {…}{{...}}1 {«}{{\guillemotleft}}1 {»}{{\guillemotright}}1}

% Shortcuts
\newcommand{\wac}{\textsc{wac}}
\newcommand{\opt}[1]{\texttt{-{}-#1}}
\newcommand{\champ}[1]{\texttt{#1}}
\newcommand{\fichier}[1]{\texttt{#1}}

% ============================================================ COVER PAGE
% The cover images are 2:3 (1024 x 1536), taller than an A4 page once widened
% to it: centred, they lost 9 mm at the top — the CEC logo — and 9 mm at the
% bottom. Raised 4 mm above the top of the page instead: the logo, which starts
% 6 mm below the image's edge, stays whole, and the bottom row of text keeps
% 7 mm to the page's edge; only a decorative corner below it is cut.
% Usage: \coverpage{cover_FR.png}
\newcommand{\coverpage}[1]{%
  \begin{titlepage}
  \thispagestyle{empty}
  \begin{tikzpicture}[remember picture,overlay]
    \node[inner sep=0pt,anchor=north,yshift=4mm] at (current page.north)
      {\includegraphics[width=\paperwidth]{#1}};
  \end{tikzpicture}
  \end{titlepage}}


\hypersetup{pdftitle={WAC — Documentation utilisateur},
  pdfauthor={Centre d'Excellence Cyberdéfense Aérospatiale (CEC)},
  pdfsubject={Logiciel de prélèvement d'artefacts Windows pour l'investigation numérique}}

\begin{document}

% ============================================================ COVER PAGE
\pagenumbering{roman}
\coverpage{cover_FR.png}

% ============================================================ INFORMATION PAGE
\thispagestyle{empty}
\vspace*{\fill}
{\color{marine}\Large\bfseries WAC — Windows Artefacts Collector}\par
{\color{bleu}Documentation utilisateur}\par
\medskip
\begin{tabular}{@{}ll@{}}
Version du logiciel & 1.4.0 \\
Édition du document & septembre 2026 \\
Éditeur & Centre d'Excellence Cyberdéfense Aérospatiale (CEC) \\
        & École de l'air et de l'espace, base aérienne 701, Salon-de-Provence \\
\end{tabular}

\bigskip
{\small\color{gris}Ce document décrit l'emploi de \wac{} : préparation, exécution
d'une collecte, lecture et vérification des résultats, et traces laissées sur la
machine examinée. Les mesures citées ont été relevées sur une machine virtuelle
Windows~11 de test ; elles donnent des ordres de grandeur, pas des
engagements.}
\vspace*{2cm}
\clearpage

% Numbering starts at 1 with the table of contents; only the cover page and the
% information page, which have no footer, are left out of it.
\pagenumbering{arabic}
\tableofcontents

% ============================================================ 1
\chapter{Présentation}

\section{À quoi sert WAC}

\wac{} prélève, sur un poste Windows en fonctionnement, les artefacts utiles à
une investigation numérique : base de registre, journaux d'événements, fichiers
Prefetch, raccourcis et listes de sauts, tâches planifiées, services, comptes,
processus et sessions. Il les restitue sous forme de fichiers \textbf{JSON},
prêts à être analysés ou versés dans une chaîne de traitement.

Il s'exécute depuis une clé USB, sans installation, et écrit tout ce qu'il
produit sur cette clé.

\section{Le principe directeur : une empreinte minimale}

Examiner une machine la modifie. Chaque outil lancé laisse des traces : fichiers
ouverts dont la date de dernier accès change, services sollicités qui écrivent
dans les journaux, clés de registre lues par l'API. Ces traces se mêlent à celles
que l'on cherche.

\wac{} est conçu pour que ce qu'il laisse derrière lui soit aussi proche de rien
que la tâche le permet, et pour que \textbf{ce qu'il ne peut pas éviter soit
écrit}. Concrètement :

\begin{itemize}
  \item le disque est lu \textbf{directement}, en lecture seule, sans ouvrir aucun
        fichier par le système : aucune date d'accès n'est modifiée ;
  \item aucune clé du registre de la machine n'est ouverte : les ruches sont
        copiées, puis lues sur la copie ;
  \item aucun cliché instantané (VSS), aucun appel COM ou WMI, aucun service de
        journalisation n'est sollicité ;
  \item rien n'est écrit sur le disque examiné ;
  \item chaque opération est consignée, avec la trace qu'elle est censée laisser,
        dans \fichier{investigation.json}.
\end{itemize}

\section{Ce que WAC ne fait pas}

\begin{itemize}
  \item \wac{} \textbf{n'analyse pas} : il prélève et met en forme. L'interprétation
        revient à l'analyste.
  \item Il ne réalise pas d'image complète du disque : il extrait les fichiers
        porteurs d'artefacts.
  \item Il ne capture pas la mémoire vive. Les seules données volatiles relevées
        sont la liste des processus, les sessions ouvertes et l'état des
        services.
\end{itemize}

% ============================================================ 2
\chapter{Principes de fonctionnement}

\section{Lecture directe du disque}

\wac{} ouvre le volume (\path|\\.\C:|) en lecture seule, parcourt lui-même la
table des fichiers NTFS (\path|$MFT|) et en extrait le contenu des fichiers
utiles. Le système ne voit aucune ouverture de fichier : ni date d'accès mise à
jour, ni parcours de répertoire, ni verrou.

Cette lecture prend en charge les particularités de NTFS que Windows masque
d'ordinaire :

\begin{description}[style=nextline,leftmargin=1.2em]
  \item[Compression NTFS] Windows~11 compresse notamment le dossier des journaux
        d'événements. \wac{} les décompresse lui-même.
  \item[Compact OS (WOF)] Les binaires système de Windows~10 et~11 sont stockés
        compressés dans un flux caché ; vus par l'API, ils paraissent ordinaires.
        \wac{} les décompresse (variantes XPRESS).
  \item[Attributs éclatés] Un fichier très fragmenté voit ses informations
        réparties sur plusieurs enregistrements ; \wac{} les réunit.
  \item[Fichiers préalloués] Au-delà de la longueur de données valides, NTFS
        rend des zéros, même si le disque contient encore les restes d'anciens
        fichiers. \wac{} fait de même : le contenu extrait est celui que verrait
        toute acquisition ordinaire, et son empreinte aussi.
\end{description}

\section{Consigne et travail}

Une pièce à conviction numérique ne s'examine jamais sur elle-même : on en fait
une copie, on la scelle, et l'on travaille sur une seconde copie. Si l'analyse
abîme quelque chose, la copie scellée demeure et l'opération peut être refaite.

\wac{} applique cette procédure à \textbf{tous} les fichiers extraits, y compris
ceux qu'il ne modifie pas :

\begin{description}[style=nextline,leftmargin=1.2em]
  \item[\fichier{exhibits/}] les copies brutes, telles que lues sur le volume. Ce
        répertoire n'est plus jamais rouvert en écriture après l'extraction. Il
        contient le manifeste qui identifie chaque pièce, et le sceau du
        manifeste.
  \item[\fichier{working/}] les copies de travail, recopiées depuis la consigne
        et vérifiées par empreinte. C'est ici que sont faites les seules
        modifications (rejeu des journaux de transaction des ruches), et c'est ici
        que lisent tous les collecteurs.
\end{description}

\begin{note}
Chaque fichier extrait est donc écrit \textbf{deux fois} sur le support de
collecte. Prévoyez la place correspondante (voir
section~\ref{sec:place}).
\end{note}

\section{Ruches du registre et journaux de transaction}

Une ruche copiée sur une machine en fonctionnement ne contient pas les dernières
modifications : elles sont encore dans ses journaux de transaction
(\fichier{.LOG1}, \fichier{.LOG2}). \wac{} extrait ces journaux et les
\textbf{rejoue sur la copie de travail}, pour que l'analyse porte sur l'état réel
de la machine.

Tout ce que \wac{} écrit sur une copie est réversible : avant le rejeu, le contenu
d'origine de chaque page remplacée est sauvegardé dans un journal d'annulation
(\fichier{<ruche>.undo}). La copie brute, elle, reste intacte dans la consigne.

\section{L'extraction en deux temps}

L'emplacement des ruches de chaque utilisateur est inscrit dans la ruche
\fichier{SOFTWARE}. \wac{} extrait donc d'abord les ruches de la machine
(\fichier{SYSTEM}, \fichier{SOFTWARE}, \fichier{SAM}, \fichier{Amcache}), lit la
liste des profils dans la copie de \fichier{SOFTWARE}, puis extrait les ruches de
chaque profil (\fichier{ntuser.dat}, \fichier{usrClass.dat}).

Rien n'est supposé sur \fichier{C:} : le lecteur système est détecté, et un
profil situé sur un autre disque (par exemple \path|D:\Users\...|) est lu sur son
propre volume.

\section{Ce qui reste lu « à chaud »}

Quelques informations décrivent l'instant même de la collecte et ne figurent
dans aucun fichier. Elles sont lues en direct, en une seule requête chacune, et
enregistrées telles que lues dans la consigne (\fichier{exhibits\textbackslash
live}), où la conversion les relit :

\begin{itemize}
  \item la liste des \textbf{processus} en cours (instantané du noyau) ;
  \item les \textbf{sessions} ouvertes ;
  \item l'\textbf{état courant des services} (la configuration, elle, est lue dans
        la ruche copiée) ;
  \item l'\textbf{horloge} et l'heure de démarrage ;
  \item les \textbf{réglages} de la machine en marche : lecteur système, page de
        code, fuseau horaire, volumes et leurs numéros de série.
\end{itemize}

% ============================================================ 3
\chapter{Préparer une collecte}

\section{Prérequis}

\begin{itemize}
  \item Windows~10 ou ultérieur, 64~bits. \wac{} le vérifie au lancement.
  \item Des \textbf{droits d'administrateur} : la lecture directe d'un volume les
        exige. \wac{} le vérifie également.
  \item Des volumes \textbf{NTFS}, \textbf{FAT} ou \textbf{exFAT}. Un fichier
        situé sur un autre système de fichiers (CD-ROM) ne peut pas être lu
        directement ; il est signalé comme non lu, jamais ouvert autrement.
\end{itemize}

\section{Le support de collecte}\label{sec:place}

\wac{} s'exécute depuis le support de collecte (en général une clé USB) et y écrit
tous ses résultats. Choisissez un support \textbf{vide}, de préférence en USB~3 :
l'écriture est le poste le plus coûteux d'une collecte.

La place nécessaire dépend du mode et des options. Deux emplacements se
remplissent : la \textbf{consigne} (copies brutes scellées) et la \textbf{copie
de travail}, qui ne reprend que ce que les collecteurs modifient ou lisent
pendant la collecte (ruches, Prefetch, raccourcis, listes de sauts, binaires
cités). Les binaires prélevés par \opt{collect} \opt{binary} et les
métadonnées NTFS ne sont jamais recopiés : la conversion les lit dans la
consigne même. Mesures sur la machine virtuelle de test (Windows~11 sans
logiciel ajouté, 1~volume) :

\begin{center}
\begin{tabular}{@{}L{7.4cm}rrr@{}}
\toprule
Mode & Consigne & Copie de travail & Durée \\
\midrule
Exécution complète, \opt{events} \opt{binary} (binaires cités) & 1,2~Go & 0,6~Go & 3~min~30 \\
Collecte seule, \opt{collect} \opt{events} \opt{binary} (tous les exécutables) & 2,7~Go & 0,3~Go & 8~min~30 \\
\bottomrule
\end{tabular}
\end{center}

Les postes les plus lourds de la consigne :

\begin{center}
\begin{tabular}{@{}L{9.6cm}r@{}}
\toprule
Contenu & Consigne \\
\midrule
Métadonnées NTFS : \path|$MFT| (380~Mo), \path|$LogFile| (64~Mo),
  \path|$UsnJrnl:$J| (35~Mo, sur 3,6~Go logiques : NTFS a libéré le début,
  qui n'est pas copié) & 480~Mo \\
Binaires copiés, exécution complète : les 62 binaires cités non authentifiés & 250~Mo \\
Binaires copiés, collecte seule : 2\,853 non authentifiés sur 44\,446 lus & 1,2~Go \\
Journaux d'événements (\opt{events}) & 170~Mo \\
Ruches, Prefetch, raccourcis, listes de sauts & 180 à 230~Mo \\
\bottomrule
\end{tabular}
\end{center}

Une machine utilisée depuis longtemps, ou chargée de logiciels, dépasse ces
chiffres : la \path|$MFT| croît avec le nombre de fichiers, les binaires tiers
non authentifiés avec les logiciels installés. Prévoir \textbf{au moins 5~Go
libres} pour une collecte seule avec \opt{binary}.

\wac{} vérifie l'emplacement \textbf{avant la toute première écriture} et refuse
de démarrer :
\begin{itemize}
  \item si le répertoire de travail contient déjà des fichiers — une collecte
        antérieure serait analysée à la place de celle-ci ;
  \item si la place libre est manifestement insuffisante.
\end{itemize}
Si la place vient à manquer en cours de prélèvement des binaires, ceux-ci sont
hachés sans être copiés, et le journal d'investigation l'indique.

\section{Avant d'intervenir}

\begin{itemize}
  \item Utilisez un \textbf{dossier de sortie neuf} pour chaque machine
        (option \opt{output}).
  \item Notez l'heure de l'intervention et l'identité de l'opérateur : \wac{} les
        consigne aussi, mais une trace indépendante est utile.
  \item Gardez à l'esprit que brancher la clé et lancer \wac{} laissent des traces
        inévitables (chapitre~\ref{chap:traces}) : elles sont connues et
        documentées, afin de ne pas être confondues avec l'activité étudiée.
\end{itemize}

% ============================================================ 4
\chapter{Lancer une collecte}

\section{Ligne de commande}

Ouvrez une invite de commandes \textbf{en tant qu'administrateur}, placez-vous
sur le support de collecte, puis lancez :

\begin{lstlisting}
WAC.exe [--collect | --convert=dossier] [--events] [--binary | --binary-all]
        [--threads=N] [--dump]
        [--output=dossier] [--loglevel=N] [--debug]
WAC.exe --update-trust[=dossier]
\end{lstlisting}

Toutes les options sont facultatives et désactivées par défaut.

\section{Fichier de configuration : \fichier{wac.yml}}

\textbf{Version.} \wac{} numérote ses versions \texttt{<majeur>.<mineure>.<correctif>} :
majeur pour une refonte, mineure pour un ajout de fonctionnalités, correctif
pour des corrections. La version se lit dans les propriétés de
\fichier{WAC.exe} (onglet Détails), s'affiche au lancement, et est consignée
dans \fichier{investigation.json} (\champ{Tool.Version}) : une collecte dit par
quelle version elle a été faite. Ce document décrit la version 1.4.0.

La dernière version compilée se trouve dans \fichier{bin} du dépôt :
\fichier{WAC.exe}, le \fichier{wac.yml} de référence, les scripts qui
vérifient une collecte (section~\ref{sec:verif}) et \fichier{SHA256SUMS}
pour vérifier le tout. Copier \fichier{WAC.exe} et \fichier{wac.yml} sur la
clé de collecte suffit ; les scripts restent avec l'enquêteur.

Plutôt que de retaper les options à chaque collecte, la procédure s'écrit une
fois dans \fichier{wac.yml}, à côté de \fichier{WAC.exe} : \wac{} l'applique
d'abord, et une option donnée en ligne de commande la surcharge.
\opt{write-config} écrit le fichier de référence, commenté — la procédure
recommandée : collecte seule, binaires non vérifiés prélevés, journaux
d'événements compris.

\begin{lstlisting}
convert: false          # true : collecte et conversion ici ; false : collecte seule
binary: unverified      # none | unverified | all
threads: 0              # fils d'analyse des binaires, 0 : automatique
output: output
log_level: 0
dump: false
language: fr            # fr | en : langue de la console ; les fichiers restent en anglais
artefacts:              # un interrupteur par artefact ; toutes les ruches = registry
  registry: true
  events: true
  scheduled_tasks: true
  services: true
  processes: true
  sessions: true
  network: true
  prefetch: true
  jump_lists: true
  recent_documents: true
  mft: true             # $MFT : chaque fichier, présent ou supprimé, et ses dates
  usn_journal: true     # $UsnJrnl:$J : créations, renommages, suppressions
  log_file: true        # $LogFile : journal de transactions NTFS, copie brute seule
  powershell_history: true  # PSReadLine : les commandes tapées dans PowerShell
  setupapi: true        # setupapi.dev.log : installations de périphériques
  recycle_bin: true     # $I des corbeilles : quoi, par qui, quand
  hosts: true           # le fichier hosts : noms détournés vers une adresse
  error_reports: true   # rapports d'erreurs (WER) : programmes plantés, avec leur chemin
  defender: true        # Defender : journaux (MPLog), historique des détections, quarantaine (copiée chiffrée)
  startup_folders: true # dossiers de démarrage : ce qui s'exécute à l'ouverture de session
  zone_identifiers: true  # flux Zone.Identifier : d'où viennent les fichiers téléchargés
  security: true        # ruche SECURITY : noms des secrets LSA, connexions en cache, audit
  persistence: true     # persistance du registre : IFEO, AppInit, Winlogon, LSA, COM…
  office_trust_records: true  # documents Office approuvés, macros activées
  typed_paths: true     # chemins tapés dans l'Explorateur
  word_wheel_query: true  # recherches de l'Explorateur
  root_certificates: true # certificats racines de la machine et des utilisateurs
  internet_settings: true # proxy des utilisateurs
  firewall: true        # règles et profils du pare-feu
  wmi: true             # dépôt WMI lu hors ligne : abonnements (filtres, consommateurs, liaisons)
  srum: true            # SRUDB.dat : réseau et ressources par application et utilisateur, 30 à 60 jours
  bits: true            # qmgr.db : travaux BITS, URL, destinations, commandes de notification
  fat_volumes: true     # volumes FAT et exFAT, clés USB du suspect comprises : entrées, supprimées aussi
  installed_programs: true  # programmes installés (Uninstall) de la machine et de chaque utilisateur
  defender_configuration: true  # exclusions et réglages de protection de Microsoft Defender
  explorer_recent_docs: true    # fichiers ouverts par chaque utilisateur (RecentDocs de l'Explorateur)
  mount_points: true            # volumes et partages vus par chaque utilisateur (MountPoints2)
  known_networks: true          # réseaux auxquels la machine s'est connectée (NetworkList)
  program_usage: true           # programmes lancés ou utilisés (FeatureUsage, JumplistData, compatibilité, épinglés)
  devices: true                 # périphériques USB et SCSI et leurs dates, portables, noms de volumes
  network_shares: true          # partages réseau de la machine, chemin et droits
  network_configuration: true   # cartes réseau, dernière configuration IPv4/IPv6, baux DHCP (hors ligne)
  event_log_configuration: true # état de chaque journal et session ETW de démarrage : actif, taille, rétention, accès (hors ligne)
\end{lstlisting}

\begin{itemize}
  \item \textbf{Strict} : une clé inconnue, une valeur hors de celles admises
    (\texttt{true}/\texttt{false} seulement, pas \texttt{yes}), une clé en
    double, une liste : la collecte est refusée avant de commencer, avec la
    ligne fautive. Une clé absente garde sa valeur par défaut.
  \item \textbf{Consigné} : chemin, SHA-256 et contenu du fichier vont dans
    \fichier{investigation.json} (\champ{Tool.Configuration}) — la procédure
    appliquée fait partie de la preuve.
  \item Un artefact désactivé n'est pas lu du tout, et son fichier JSON le dit
    (\champ{CollectionStatus} \texttt{NotRequested}, avec la clé) : jamais
    confondu avec un échec ou une absence de traces.
  \item \opt{config=fichier} lit un autre fichier ; \opt{full} collecte et
    convertit ici bien que le fichier dise \texttt{convert: false}.
  \item Sans \fichier{wac.yml}, \wac{} se comporte comme avant : exécution
    complète, options telles que données.
\end{itemize}

\section{Options}

\begin{longtable}{@{}>{\ttfamily}L{3.2cm}L{11.2cm}@{}}
\toprule
\normalfont\bfseries Option & \bfseries Effet \\
\midrule
\endhead
-{}-collect & Collecte seule, sur la machine examinée : relevés en direct et
  extraction brute dans la consigne scellée, rien n'est converti. Voir
  chapitre~\ref{chap:separation}. \\
-{}-convert=\newline dossier & Conversion seule, sur un poste d'analyse, de la collecte
  faite avec \opt{collect} dans ce dossier, après vérification du sceau et de
  chaque empreinte. Voir chapitre~\ref{chap:separation}. \\
-{}-events & Extrait et analyse les journaux d'événements (\fichier{.evtx}), et
  reconstitue le message lisible de chaque événement à partir des ressources de
  son fournisseur. Ajoute environ 250~Mo de copies. \\
-{}-binary & Calcule les empreintes (MD5, SHA-1, SHA-256) de chaque fichier cité
  par un artefact, par lecture directe du disque. Les exécutables, bibliothèques,
  pilotes et scripts cités sont en outre \textbf{prélevés} dans la consigne. Voir
  section~\ref{sec:binary}. \\
-{}-binary-all & Comme \opt{binary}, mais \textbf{tout} exécutable est prélevé,
  authentifié ou non ; la signature est tout de même vérifiée, et le manifeste
  dit pour chacun si elle est valide. Pour les dossiers où le contenu même doit
  figurer dans la consigne. \\
-{}-threads=N & Nombre de fils qui analysent les exécutables de \opt{collect}
  \opt{binary} (PE, Authenticode, catalogues, empreintes). Par défaut, le nombre
  de fils du processeur moins un, la lecture du disque gardant le sien. Le
  résultat n'en dépend pas : le manifeste est le même. \\
-{}-update-trust\newline [=dossier] & Sur le \textbf{poste d'analyse}, \textbf{avant} la collecte, et
  seule : télécharge chez Microsoft la liste des racines de confiance
  (\fichier{authroot.stl}), celle des certificats désapprouvés
  (\fichier{disallowedcert.stl}) et chacune des racines, vérifie leurs
  signatures et leurs empreintes, ainsi que les listes de pilotes vulnérables
  de Microsoft et de LOLDrivers (non signées : authentifiées par HTTPS seul)
  et les listes de révocation (CRL) des autorités de signature de code que
  recense la CCADB, de celles propres à Microsoft et de celles qu'une collecte
  précédente a trouvées manquantes (\fichier{wanted-crls.txt}), et les écrit dans \fichier{dossier} (par
  défaut \fichier{trust}, à côté de \fichier{WAC.exe}), à emporter sur la clé
  avec \fichier{WAC.exe}. Seul mode qui utilise le réseau : la bibliothèque
  réseau n'est chargée que là. Refusée avec toute autre option. \\
-{}-dump & Ajoute la valeur hexadécimale brute des éléments de shellbags et de
  raccourcis \fichier{.lnk}, pour un décodage ultérieur. \\
-{}-output=dossier & Dossier de sortie, relatif au répertoire courant. Par défaut :
  \fichier{output}. \\
-{}-language=fr|en & Langue de la console : étapes, aide, messages d'erreur.
  Français par défaut ; \texttt{language:} dans \fichier{wac.yml} la fixe
  aussi. Les fichiers JSON, \fichier{investigation.json} et le journal texte
  restent en anglais, quelle que soit la langue : ce sont des pièces lues par
  des outils, et une collecte donne les mêmes fichiers quel que soit
  l'opérateur. \\
-{}-loglevel=N & Active un journal texte détaillé, dans le répertoire courant,
  nommé d'après le programme lancé (\fichier{WAC.exe.log} pour \fichier{WAC.exe}) : 0~aucun, 1~par type
  d'artefact, 2~par artefact, 3~par sous-fonction (mise au point). \\
-{}-debug & Affiche sur la sortie d'erreur la trace du lecteur NTFS (résolution
  des chemins, blocs d'index, séquences de grappes). Réservé au diagnostic. \\
-{}-help, /? & Affiche l'aide. \\
\bottomrule
\end{longtable}

\section{Exemples}

Collecte standard :
\begin{lstlisting}
E:\> WAC.exe --output=poste-comptable-01
\end{lstlisting}

Collecte complète, avec journaux, binaires et journal détaillé :
\begin{lstlisting}
E:\> WAC.exe --events --binary --output=poste-comptable-01 --loglevel=2
\end{lstlisting}

\section{L'option \opt{binary}}\label{sec:binary}

Une empreinte suffit pour interroger une base publique (VirusTotal, par exemple)
\textbf{sans lui envoyer le fichier}. Mais elle ne dit rien d'un binaire que
personne ne connaît — précisément celui qui intéresse l'enquête — et un binaire
non prélevé peut avoir disparu quand une détection survient. C'est pourquoi
\opt{binary} copie aussi les charges potentielles :

\begin{itemize}
  \item \textbf{prélevés} : \fichier{exe}, \fichier{dll}, \fichier{sys},
        \fichier{ocx}, \fichier{cpl}, \fichier{scr}, \fichier{drv},
        \fichier{efi}, \fichier{com}, \fichier{msi} et les scripts
        \fichier{ps1}, \fichier{psm1}, \fichier{bat}, \fichier{cmd},
        \fichier{vbs}, \fichier{vbe}, \fichier{js}, \fichier{jse},
        \fichier{wsf}, \fichier{wsh}, \fichier{hta} ;
  \item \textbf{prélevés aussi} : les documents Office \textbf{capables de porter
        des macros}, vecteur d'intrusion aussi courant qu'un exécutable —
        formats binaires anciens (\fichier{doc}, \fichier{xls}, \fichier{ppt},
        \fichier{pps}, \fichier{dot}, \fichier{xlt}, \fichier{pot}), formats
        « m » (\fichier{docm}, \fichier{dotm}, \fichier{xlsm}, \fichier{xltm},
        \fichier{pptm}, \fichier{potm}, \fichier{ppsm}), \fichier{xlsb},
        compléments (\fichier{xla}, \fichier{xlam}, \fichier{xll},
        \fichier{wll}, \fichier{ppa}, \fichier{ppam}), Publisher
        (\fichier{pub}), Visio (\fichier{vsd}, \fichier{vsdm}, \fichier{vstm},
        \fichier{vssm}) et Access (\fichier{mdb}, \fichier{accdb},
        \fichier{accde}). Ils apparaissent dans les traces comme cibles de
        raccourcis ou de listes de sauts, ou parmi les fichiers chargés par
        Word ou Excel dans leur Prefetch ;
  \item \textbf{seulement hachés} : les autres documents et fichiers de données,
        y compris \fichier{docx}, \fichier{xlsx} et \fichier{pptx}, qui ne
        peuvent pas contenir de macros. Les copier ferait de la collecte une
        copie des fichiers de l'utilisateur.
\end{itemize}

Les fichiers cités proviennent des processus, services (binaire et DLL), tâches
planifiées, fichiers chargés par les programmes (Prefetch), Shimcache, Amcache et
cibles de raccourcis et de listes de sauts. Chaque fichier n'est lu qu'une fois, quel que soit le
nombre d'artefacts qui le citent, et \textbf{un même contenu n'est stocké
qu'une fois} : les autres chemins figurent au manifeste comme pièces partageant
ce contenu (\champ{SharedExhibit}).

\subsection*{Binaires Microsoft authentiques : hachés, non prélevés}

La plupart des binaires cités sont des composants de Windows, identiques sur
toute machine de la même version. Leur origine se prouve sur la machine
elle-même, sans liste à embarquer : les \textbf{catalogues} de Windows, signés
par Microsoft, listent l'empreinte de chaque fichier du système. \wac{} lit ceux
de \path|System32\CatRoot|, ceux des composants (\path|WinSxS\Catalogs|) et
ceux des paquets de maintenance (\path|servicing\Packages|), plus de 13\,000 sur
Windows~11. Les binaires signés individuellement (Edge, OneDrive, Office, .NET)
portent une \textbf{signature intégrée}. Un binaire authentifié ainsi est
haché sans être prélevé ; tous les autres sont prélevés — y compris un binaire
système remplacé par un attaquant, dont l'empreinte ne correspond plus.

Cette vérification ne laisse \textbf{aucune trace} : catalogues et binaires sont
lus directement sur le disque, et la vérification cryptographique est faite en
mémoire par \wac{}, jusqu'à des certificats racines Microsoft embarqués dans
l'exécutable. Aucun service n'est sollicité, le magasin de certificats de la
machine n'est pas consulté, et la révocation n'est pas contrôlée en ligne.

Seuls les signataires Microsoft pour son propre code sont acceptés
(« Microsoft Windows », « Microsoft Corporation », « Microsoft Windows
Publisher » et « .NET »). Les pilotes
\textbf{tiers} certifiés par Microsoft (signataires « Hardware Compatibility
Publisher » et « Early Launch Anti-malware Publisher ») restent prélevés : un
pilote vulnérable mais signé est un moyen d'attaque classique.

\subsection*{Applications du Store : vérifiées par leur paquet}

Les fichiers d'une application du Store ne sont pas signés un à un : c'est leur
\textbf{paquet} qui l'est. Son dossier contient la signature
(\fichier{AppxSignature.p7x}) et la carte des blocs (\fichier{AppxBlockMap.xml}),
qui donne le SHA-256 de chaque bloc de 64~Kio de chaque fichier. \wac{} vérifie
la signature, qui doit remonter à une racine Microsoft (le Store signe les
paquets qu'il distribue), vérifie qu'elle garantit bien cette carte, puis compare
chaque bloc lu à la carte. Un fichier conforme est haché sans être prélevé ; un
fichier modifié, ou ajouté au dossier après l'installation, est prélevé, de même
que tout fichier d'un paquet signé hors de la chaîne Microsoft.

Les catalogues, signatures et cartes de blocs qui ont justifié un
non-prélèvement sont consignés, pour qu'un tiers puisse refaire la vérification.
Dans les fichiers d'artefacts, le champ \champ{Signature} (ou
\champ{ServiceDllSignature}, \champ{SignatureTarget}\ldots) indique
l'authentification retenue, par exemple
« \texttt{Microsoft (catalog C:\textbackslash{}Windows\textbackslash{}\ldots.cat)} »,
« \texttt{Microsoft (embedded signature)} » ou
« \texttt{Package (package C:\textbackslash{}Program Files\textbackslash{}WindowsApps\textbackslash{}\ldots)} ».

Les \textbf{scripts} sont authentifiés de la même façon : par les catalogues de
Windows, et, pour PowerShell (\fichier{ps1}, \fichier{psm1}, \fichier{psd1},
\fichier{ps1xml}\ldots), par leur bloc de signature intégré
(\texttt{\# SIG \# Begin signature block}). Un script Windows Script Host
(\fichier{vbs}, \fichier{js}, \fichier{wsf}) signé seulement dans son texte, et
non par un catalogue, reste prélevé.

Sur la machine de test, cette vérification ramène le prélèvement de 2~131
binaires (3,3~Go) à 92 (335~Mo), sans allonger la collecte ; les scripts encore
prélevés sont ceux qui ne sont pas signés.

La lecture étant directe, prélever ne laisse aucune trace de plus que hacher :
seulement de la place sur le support.

\section{Déroulement à l'écran}

\wac{} affiche chaque étape, suivie de \texttt{OK} ou d'un message d'erreur. Une
étape en échec n'interrompt pas la collecte : tout ce qui est accessible est
recueilli, et ce qui manque est consigné.

\begin{lstlisting}
[VÉRIFICATION DES PRÉREQUIS]
 - Windows 10 ou ultérieur : OK
 - Droits d'administrateur : OK
 - Contrôle du support de collecte : OK
[OBSERVATION DE LA MACHINE EN MARCHE]
 - Sessions : OK
 - Processus : OK
 - État du réseau : OK
[EXTRACTION BRUTE]
 - Ruches du système (lecture brute NTFS) : OK
 - Ouverture de la ruche SYSTEM : OK
 - Ouverture de la ruche SOFTWARE : OK
 - Profils des utilisateurs (ruche SOFTWARE) : OK
 - Ruches des utilisateurs (lecture brute NTFS) : OK
 - Artefacts fichiers (lecture brute NTFS) : OK
[ARTEFACTS DU REGISTRE]
 - Clés UserAssist du registre : OK
   ...
[ARTEFACTS DES FICHIERS]
 - Prefetch : OK
   ...
[JOURNAUX D'ÉVÉNEMENTS]
 - Journaux d'événements : OK
[CONSIGNE]
 - Copie des binaires prélevés dans la copie de travail : OK
 - Scellement de la consigne (manifeste + SHA-256) : OK
[JOURNAL D'INVESTIGATION]
 - Écriture de investigation.json : OK
FIN, durée : 309 s
\end{lstlisting}

Avec \opt{language=en}, ou \texttt{language: en} dans \fichier{wac.yml}, les
mêmes étapes s'affichent en anglais.

Le scellement de la consigne est la \textbf{dernière} opération sur celle-ci :
toutes les pièces, y compris celles ajoutées pendant l'analyse des journaux, sont
couvertes par le manifeste.

\section{Durée}

La durée dépend surtout du matériel (USB~2 ou~3, disque, processeur) et des
options. Sur la machine virtuelle de test (disque NVMe) : environ 70~s avec
\opt{events}, environ 310~s avec \opt{events} et \opt{binary}. Sur une clé USB,
l'écriture domine.

% ============================================================ 4 bis
\chapter{Collecter ici, convertir ailleurs}\label{chap:separation}

\section{Pourquoi séparer}

Une exécution de \wac{} fait deux choses différentes :
\begin{enumerate}
  \item \textbf{la collecte} : relever ce que seul un système en marche montre
        (processus, sessions, état des services, horloge), lire le disque
        directement, copier les preuves dans une consigne scellée. Elle a
        \textbf{besoin} de la machine examinée ;
  \item \textbf{la conversion} : décoder en JSON des milliers de fichiers
        copiés (ruches, journaux, Prefetch, raccourcis\ldots). Elle n'a besoin
        que de la consigne, \textbf{pas} de la machine.
\end{enumerate}

Sans option de mode, les deux se font sur la machine examinée. La conversion y
ajoute un travail que rien n'impose : des minutes de traitement dans la mémoire
de \wac{}, que Windows peut paginer dans \fichier{pagefile.sys}, sur le disque
même qu'on examine ; un processus présent plus longtemps, vu et journalisé par
les protections de la machine ; en cas de plantage, un rapport d'erreur Windows
écrit sur la machine. Elle ne lit pourtant ni les fichiers ni le registre de la
machine, seulement les copies du support de collecte, mais elle s'y exécute.

C'est précisément ce que les bonnes pratiques demandent d'éviter :
\begin{itemize}
  \item \textbf{ISO/IEC~27037} traite de l'identification, de la collecte et de
        la préservation des preuves numériques, \textbf{ISO/IEC~27042} de leur
        analyse : deux étapes distinctes ;
  \item \textbf{principe~1 de l'ACPO} : aucune action ne doit modifier les
        données de la machine examinée ; toute opération qui s'y exécute doit
        être nécessaire ;
  \item \textbf{principe~3 de l'ACPO} : un tiers indépendant doit pouvoir
        refaire le traitement et obtenir le même résultat.
\end{itemize}

Séparer les deux étapes apporte trois choses :
\begin{itemize}
  \item \textbf{moins de traces} sur la machine examinée : seule la collecte
        s'y exécute ;
  \item \textbf{des preuves vérifiées avant toute analyse} : \opt{convert}
        contrôle d'abord le sceau et le SHA-256 de chaque pièce, et refuse une
        collecte retouchée ;
  \item \textbf{une analyse reproductible} : deux conversions d'une même
        collecte donnent des fichiers JSON identiques octet pour octet, et
        l'analyse peut être refaite plus tard, avec une version corrigée de
        \wac{} ou par un tiers, sans revenir sur la machine.
\end{itemize}

\section{Quel mode choisir}

\begin{itemize}
  \item \textbf{Par défaut : \opt{collect} sur la machine, \opt{convert} sur le
        poste d'analyse.} C'est la procédure qui laisse le moins de traces sur la
        machine examinée, garde les preuves scellées et permet de refaire et de
        vérifier l'analyse.
  \item \textbf{Exécution complète} (sans option de mode) : à réserver aux cas
        particuliers, à signaler au rapport — un tri dont les résultats sont
        attendus sur place sans poste d'analyse, une machine à rendre en quelques
        minutes. Sa conversion ne lit que les copies du support de collecte,
        jamais les fichiers ni le registre de la machine, mais elle s'exécute sur
        celle-ci : des minutes de traitement en mémoire (que Windows peut paginer
        dans \fichier{pagefile.sys}), un processus plus long, un rapport
        d'erreur Windows en cas de plantage.
  \item \textbf{Portée de \opt{binary}} : en exécution complète, seuls les
        fichiers \textbf{cités par les artefacts} sont lus ; un binaire du disque
        qu'aucun artefact ne cite n'est pas examiné. Avec \opt{collect}, ce sont
        \textbf{tous les exécutables de tous les volumes fixes NTFS}, cités ou non.
        La collecte seule dure plus longtemps pour cette raison (8~min~30
        contre 3~min~30 sur la machine de test) : ce temps est de la lecture,
        qui n'écrit rien sur la machine.
  \item \textbf{\opt{binary-all}} : quand le contenu de chaque exécutable doit
        figurer dans la consigne (analyse de logiciel malveillant, demande des
        fichiers eux-mêmes), pour environ 20~Go.
\end{itemize}

\section{Les deux commandes}

Sur la machine examinée, en administrateur :
\begin{lstlisting}
E:\> WAC.exe --collect --events --binary --output=poste-comptable-01
\end{lstlisting}

Sur le poste d'analyse, sans droits particuliers, avec les mêmes options
\opt{events} et \opt{binary} que la collecte :
\begin{lstlisting}
D:\> WAC.exe --convert=E:\poste-comptable-01 --events --binary
\end{lstlisting}

Les fichiers JSON sont écrits dans le dossier de la collecte, à côté de la
consigne. \opt{convert} exclut \opt{collect} et \opt{output}.

\section{Ce que prend \opt{collect}}

La conversion ne pouvant plus rien demander à la machine, la collecte prend
\textbf{tout ce qu'elle lira} :
\begin{itemize}
  \item les \textbf{relevés en direct}, enregistrés tels que lus dans
        \fichier{exhibits\textbackslash live} et scellés comme toute pièce :
        processus, sessions, état des services, horloge et démarrage, et les
        réglages de la machine (lecteur système, page de code, fuseau horaire,
        volumes et leurs numéros de série) ;
  \item les ruches et les fichiers d'artefacts, comme d'ordinaire ;
  \item avec \opt{events}, les fichiers de ressources de \textbf{tous} les
        fournisseurs d'événements déclarés (binaire, fichier de paramètres,
        satellites \fichier{.mui}), car on ne sait qu'à la conversion lesquels
        seront utiles. Cela représente environ 650~Mo sur Windows~11 ;
  \item avec \opt{binary}, \textbf{tous} les exécutables, bibliothèques,
        pilotes, scripts et documents à macros de \textbf{tous les volumes fixes
        NTFS} — un profil sur \fichier{D:} compris —, lus directement et
        vérifiés en mémoire, comme en exécution complète (section~\ref{sec:binary}) :
        \begin{itemize}
          \item authentifié (catalogue, signature intégrée, paquet du Store), ou
                binaire tiers dont la signature tient face au jeu de confiance de
                la clé (\opt{update-trust} : chaîne vers une racine admise par
                Microsoft pour la signature de code, valide à la date de
                signature, non révoquée, pas un pilote vulnérable) :
                \textbf{non copié} ; sans jeu de confiance utilisable, tout
                binaire tiers est copié, et le journal d'investigation le dit ; le manifeste garde son chemin, ses empreintes
                Authenticode et le verdict (\champ{ContentStored} à \texttt{false}) ;
          \item non authentifié : \textbf{copié}, avec ses trois empreintes
                MD5, SHA-1, SHA-256, son empreinte Authenticode et la raison ;
          \item fichier différentiel de Windows (\texttt{PA30}/\texttt{PA31},
                dossiers \fichier{r} et \fichier{f} de \fichier{WinSxS}) : ce n'est
                pas un exécutable ; haché, non copié.
        \end{itemize}
        Un même contenu n'est stocké qu'une fois, et les autres noms d'un lien
        physique (\fichier{System32} et \fichier{WinSxS}) sont déclarés au
        manifeste comme le premier. Les volumes amovibles, dont le support de
        collecte, et les fichiers de mémoire (\fichier{pagefile.sys},
        \fichier{swapfile.sys}, \fichier{hiberfil.sys}) sont exclus.
\end{itemize}

\begin{attention}
Avec \opt{binary}, la collecte lit tous les volumes : sur la machine
virtuelle de test (Windows~11 sans logiciel ajouté), 44\,446~exécutables lus,
dont 41\,593 authentifiés (ou tiers à signature valide) non copiés et
2\,853 copiés (1,2~Go) ; la collecte entière dure 8~minutes~30 et la consigne,
journaux d'événements et métadonnées NTFS compris, pèse 2,7~Go (voir la
section~\ref{sec:place}). Si la place vient à manquer, les fichiers
restants sont hachés sans être copiés, et le manifeste l'indique.
\end{attention}

Aucun fichier JSON d'artefact n'est écrit : le dossier ne contient que la
consigne scellée, le répertoire de travail et \fichier{investigation.json}.

\section{Ce que fait \opt{convert}}

\begin{enumerate}
  \item \textbf{Vérification}, avant toute lecture : le sceau doit porter
        l'empreinte réelle du manifeste, et chaque pièce son empreinte SHA-256.
        Une seule pièce altérée, et la conversion est \textbf{refusée} ;
        \fichier{conversion.json} en donne la raison.
  \item \textbf{Reconstruction du répertoire de travail} depuis la consigne :
        l'ancien est supprimé plutôt que cru, les copies sont vérifiées contre
        le manifeste et les ruches réparées, comme à la collecte. Les
        exécutables et les ressources des fournisseurs sont lus dans la
        consigne même, sans copie.
  \item \textbf{Conversion}, depuis les seules pièces : aucune lecture de la
        machine qui convertit, ni de son disque, ni de son registre, ni de son
        fuseau horaire.
\end{enumerate}

La consigne n'est jamais modifiée et son sceau reste celui de la collecte.
\fichier{investigation.json} reste le journal de la collecte ;
\fichier{conversion.json} est celui de la conversion (opérateur, poste,
opérations). Deux conversions d'une même collecte donnent des fichiers JSON
identiques, ce que le harnais de test vérifie à chaque cycle.

\section{Limites}

\begin{itemize}
  \item Les \textbf{documents à macros} (\fichier{doc}, \fichier{docm},
        \fichier{xls}, \fichier{xlsm}, \fichier{ppt}\ldots, liste de la
        section~\ref{sec:binary}) sont traités comme les exécutables : lus sur
        tous les volumes fixes et \textbf{prélevés}, sauf s'ils sont, octet pour
        octet, un fichier livré par Microsoft et listé dans un catalogue — ce
        qu'un document écrit ou modifié par quiconque ne peut pas être.
  \item Un fichier cité qui ne peut porter ni code ni macro (\fichier{docx},
        \fichier{xlsx}, \fichier{pdf}, donnée) n'est pas prélevé : en
        conversion, il n'a donc \textbf{pas d'empreinte}, là où l'exécution
        complète la calcule sur le disque.
  \item Un exécutable authentifié n'est pas copié : à la conversion, il a son
        empreinte Authenticode et son verdict, mais pas de MD5, SHA-1 ni
        SHA-256 — ceux-ci ne sont calculés que pour ce qui est prélevé.
  \item Un exécutable situé sur un volume amovible n'est pas prélevé par
        \opt{collect} \opt{binary}.
  \item Les options \opt{events} et \opt{binary} de la conversion doivent être
        celles de la collecte : ce que la collecte n'a pas pris ne peut pas être
        converti.
\end{itemize}

% ============================================================ 5
\chapter{Les résultats}

\section{Arborescence}

\begin{lstlisting}
<sortie>/
  exhibits/            copies brutes, jamais réécrites
    MANIFEST.json     identification de chaque pièce et de la collecte
    MANIFEST.sha256   empreinte du manifeste : son sceau
    Windows/...        les pièces, sous leur chemin d'origine
    Users/...
  working/             copies de travail (ruches rejouées)
  investigation.json   journal de la collecte
  *.json               les artefacts
\end{lstlisting}

Un fichier d'un autre volume que le volume système est rangé sous
\path|_volume_D\...| ; le JSON indique toujours le chemin d'origine, avec sa
vraie lettre de lecteur.

\section{Le manifeste de consigne}

\fichier{exhibits/MANIFEST.json} identifie la collecte et chaque pièce.

\subsection{En-tête}

\begin{description}[style=nextline,leftmargin=1.2em]
  \item[\champ{Tool}] nom de l'outil, ligne de commande, date de compilation.
  \item[\champ{Host}] nom de la machine, lecteur système, fuseau horaire du
        suspect (lu dans sa ruche \fichier{SYSTEM}) et de la machine de collecte.
  \item[\champ{Operator}] compte ayant lancé la collecte, et s'il était élevé.
  \item[\champ{Custody}] déclaration de procédure, début de l'extraction, heure
        d'écriture du manifeste, nombre de pièces et d'échecs, volume total,
        algorithmes d'empreinte, volumes lus (lettre, numéro de série, système de
        fichiers).
\end{description}

\subsection{Chaque pièce}

\begin{longtable}{@{}>{\ttfamily}L{4.2cm}L{10.2cm}@{}}
\toprule
\normalfont\bfseries Champ & \bfseries Signification \\
\midrule
\endhead
SourcePath & Chemin d'origine, avec lettre de lecteur. \\
ExhibitPath & Emplacement de la copie brute, relatif au dossier de sortie. \\
Method & Méthode de collecte (lecture directe NTFS, volume, motif). \\
Result & \texttt{OK}, ou le code et le message d'erreur. Les échecs figurent au
  manifeste : une pièce absente se lirait comme une pièce jamais cherchée. \\
MD5, SHA1, SHA256 & Les trois empreintes, calculées pendant l'écriture : elles
  portent sur ce qui a été lu du volume. \\
Bytes & Taille extraite. \\
DeclaredBytes & Taille déclarée par NTFS, présente seulement si elle diffère de
  la taille extraite (extraction tronquée). \\
ValidDataBytes & Longueur de données valides d'un fichier préalloué (journaux
  d'événements) : au-delà, la pièce contient des zéros par définition. \\
SharedExhibit & \texttt{true} si le contenu, identique octet pour octet, est déjà
  stocké sous une autre pièce : \champ{ExhibitPath} désigne celle-ci. La source
  reste une pièce à part entière, avec ses propres chemin, entrée \path|$MFT| et
  horodatages. \\
MftEntry & Numéro d'enregistrement NTFS : identifie le fichier indépendamment
  de son nom. \\
Extracted(Utc) & Instant de l'extraction de cette pièce. \\
SourceCreated, SourceModified, SourceMftModified, SourceAccessed & Les quatre
  horodatages NTFS du fichier d'origine (\path|$STANDARD_INFORMATION|), en UTC et
  en heure locale du suspect. Ce sont eux, et non les dates des copies, qui
  datent le fichier. \\
SourceFileName\newline Created, Modified, MftModified, Accessed &
  Les quatre horodatages de l'attribut \path|$FILE_NAME| du nom lu. NTFS les
  tient à part, et une modification des dates par programme ne les touche pas :
  un \champ{SourceCreated} antérieur à \champ{SourceFileNameCreated} est le signe
  classique de dates falsifiées après coup (\emph{timestomping}). \\
ContentStored & \texttt{false} : binaire authentifié lu mais non copié
  (\opt{collect} \opt{binary}) ; il n'a pas de \champ{ExhibitPath}. \\
AuthenticodeSHA1, AuthenticodeSHA256 & Empreintes Authenticode d'un PE, celles
  que couvrent sa signature et les catalogues ; seules empreintes d'un binaire
  authentifié non copié. \\
SignatureVerified & Résultat de la vérification de signature d'un exécutable :
  \texttt{true} avec \champ{Signature} (ce qui l'authentifie), \texttt{false}
  avec \champ{SignatureReason} (pourquoi). \\
PeTimeDateStamp, PeSizeOfImage, SymbolServerKey & Version d'un PE et clé
  (\texttt{nom/HORODATAGEtaille/nom}) sous laquelle le serveur de symboles de
  Microsoft la conserve : un binaire Microsoft authentique non copié peut en
  être retéléchargé, identique, et son empreinte revérifiée. \\
\bottomrule
\end{longtable}

\begin{note}[Pourquoi trois empreintes]
Des collisions MD5 se fabriquent à volonté depuis 2008, et SHA-1 depuis 2017 : une
empreinte seule ne tranche plus une contestation d'identité. Trois empreintes
indépendantes, si.
\end{note}

\section{Vérifier l'intégrité d'une collecte}\label{sec:verif}

\subsection{Le sceau}

Le manifeste ne peut pas contenir sa propre empreinte : elle est dans
\fichier{MANIFEST.sha256}. Sous Linux :

\begin{lstlisting}
$ cd sortie/exhibits
$ sha256sum -c MANIFEST.sha256
MANIFEST.json: OK
\end{lstlisting}

\subsection{Chaque pièce}

Deux scripts, livrés avec \fichier{WAC.exe} dans \fichier{bin},
vérifient le sceau puis le SHA-256 de chaque pièce :
\fichier{verifier-consigne.ps1} sous Windows, \fichier{verifier-consigne.sh}
sous Linux. Leurs équivalents anglais, \fichier{verify-exhibits.ps1} et
\fichier{verify-exhibits.sh}, font les mêmes contrôles avec des messages en
anglais.

\textbf{Vérifier depuis une machine Linux.} C'est la voie recommandée : la
collecte peut contenir ce qui a infecté la machine examinée, et une machine
d'investigation sous Linux constitue une rupture technologique — un exécutable
ou un document à macro Windows de la collecte ne peut pas s'y exécuter. Le
script ne demande que \fichier{bash}, les coreutils et \fichier{python3} (qui
lit le JSON) ; les empreintes sont calculées par \fichier{sha256sum}, non par
\wac{} :

\begin{lstlisting}
$ bash verifier-consigne.sh /media/usb/output
Sceau conforme (DCFF53FAE709EC2E351FFFD02EC6EAD9945602265CF3363D0D928B0B452C3E32)
3360 piece(s) conforme(s), 0 differente(s)
0 binaire(s) authentifie(s), empreinte seule (pas de copie a verifier)
\end{lstlisting}

Le manifeste vient de la machine examinée : les deux scripts refusent un
chemin de pièce qui sort de la collecte (\texttt{..}, chemin absolu, et pour
le \fichier{.sh} un lien symbolique), et le comptent comme différent sans le
lire.

\textbf{Depuis Windows :}

\begin{lstlisting}
PS> powershell -ExecutionPolicy Bypass -File verifier-consigne.ps1 -Sortie E:\output
Sceau conforme (DCFF53FAE709EC2E351FFFD02EC6EAD9945602265CF3363D0D928B0B452C3E32)
3360 piece(s) conforme(s), 0 differente(s)
\end{lstlisting}

Leur code de retour vaut~0 si tout est conforme, 1~sinon.

\lstinputlisting[language={},caption={}]{verifier-consigne.sh}

\lstinputlisting[language={},caption={}]{verifier-consigne.ps1}

\section{Le journal d'investigation}

\fichier{investigation.json} retrace la collecte : outil et ligne de commande,
machine et fuseaux horaires, opérateur, début, fin et durée, puis la liste
\textbf{ordonnée} des opérations. Pour chacune :

\begin{description}[style=nextline,leftmargin=1.2em]
  \item[\champ{Seq}, \champ{TimestampUtc}, \champ{TimestampLocal}] rang et instant ;
  \item[\champ{Operation}, \champ{Target}] ce qui a été fait, et sur quoi ;
  \item[\champ{Result}] l'issue ;
  \item[\champ{Footprint}] la trace que cette opération est censée laisser sur la
        machine examinée.
\end{description}

Ce journal permet de confronter, collecte par collecte, les traces attendues à
celles effectivement observées. Il signale aussi un écart entre le fuseau horaire
du suspect et celui de la machine de collecte.

\section{Les fichiers d'artefacts}

Un fichier par artefact. Les champs décrits sont ceux émis ; un champ sans valeur
n'est pas émis (voir chapitre~\ref{chap:lecture}).

\subsection{Système et comptes}

\begin{description}[style=nextline,leftmargin=1.2em]
  \item[\fichier{OperatingSystem.json}] nom et domaine de la machine, édition et
        version de Windows, date d'installation, heure courante, \textbf{heure de
        démarrage}, durée d'activité, fuseau horaire.
        \champ{BootTimeSource} dit d'où vient l'heure de démarrage (noyau, ou
        estimation de repli) ; \champ{ClockAdjustedSinceBootMs} donne le cumul des
        corrections d'horloge depuis le démarrage — quelques secondes sont un
        recalage ordinaire, un grand écart signale une horloge modifiée.
        \champ{ShutdownTime} : le dernier arrêt ordonné.
        \champ{NtfsLastAccessUpdated} : faux, les dates d'accès NTFS ne sont pas
        mises à jour et ne disent rien. \champ{Settings} : les réglages qui
        conditionnent la lecture des autres artefacts, tels que stockés --- dates
        d'accès et noms courts NTFS, Prefetch actif, fichier d'échange et vidages
        mémoire, jeux de contrôle, bureau à distance, chronologie Windows,
        dernier utilisateur connecté.
  \item[\fichier{security.json}] ruche \fichier{SECURITY} : \textbf{noms} des
        secrets LSA et dates de leurs valeurs (\texttt{DefaultPassword} signale
        un mot de passe de connexion automatique enregistré,
        \texttt{\_SC\_<service>} un service qui tourne sous un compte),
        nombre d'emplacements de connexions de domaine en cache et combien sont
        occupés, politique d'audit (\champ{PolAdtEv}, en octets : son format
        n'est pas publié). \textbf{Aucune valeur de secret} n'est lue, ni
        aucune donnée d'authentification en cache. Les copies
        \path|config\RegBack| des ruches sont prélevées telles quelles.
  \item[\fichier{users.json}] comptes locaux, lus dans la ruche \fichier{SAM} :
        SID et RID, état, drapeaux, nombre de connexions et d'échecs, dates de
        dernière connexion, de dernier échec, de changement de mot de passe et de
        modification du compte, profil.
  \item[\fichier{Sessions.json}] sessions ouvertes au moment de la collecte :
        utilisateur, type d'ouverture, paquet d'authentification, heure de début.
  \item[\fichier{processes.json}] processus en cours : nom, identifiants (et
        parent), propriétaire, session, modules chargés, et, avec \opt{binary},
        empreintes de l'exécutable.
  \item[\fichier{network.json}] état réseau au moment de la collecte :
        \champ{Connections} (TCP et UDP, IPv4 et IPv6 : adresses, ports, état,
        processus propriétaire et son nom, heure de création locale et UTC),
        \champ{Neighbours} (caches ARP et NDP : les machines du réseau local
        contactées récemment), \champ{Interfaces} (adresses, passerelles,
        serveurs DNS, DHCP) et \champ{Routes}. \champ{TablesNotRead} nomme une
        table que Windows a refusé de donner, ce qui n'est pas une table vide.
  \item[\fichier{services.json}] services et pilotes, configuration lue dans la
        ruche : type, démarrage, compte, binaire et DLL hébergée, commande en cas
        d'échec, dépendances, dernière écriture de la clé ; état courant ; avec
        \opt{binary}, empreintes du binaire et de la DLL.
  \item[\fichier{ScheduledTasks.json}] tâches planifiées, lues dans leurs
        définitions XML : auteur, compte d'exécution, état, dernière exécution et
        son résultat, déclencheurs et actions (avec \opt{binary}, empreintes de la
        commande).
\end{description}

\subsection{Exécution de programmes}

\begin{description}[style=nextline,leftmargin=1.2em]
  \item[\fichier{prefetchs.json}] fichiers Prefetch : programme, nombre et dates
        d'exécution, volumes, et liste des fichiers chargés avec leur référence
        \path|$MFT| (et leurs empreintes avec \opt{binary}).
  \item[\fichier{amcache\_application\_files.json},
        \fichier{amcache\_applications.json},
        \fichier{amcache\_inventory.json}] inventaire Amcache, table par
        table et valeur par valeur, avec la date de dernière écriture de chaque
        entrée : les exécutables (chemin, éditeur, version, taille, date
        d'édition de liens, \champ{Sha1}), les applications installées, et dans
        le troisième fichier les pilotes (avec leur \champ{Sha1}), périphériques,
        conteneurs, paquets de pilotes et raccourcis, chacun avec sa
        \champ{Table}. \textbf{Attention} : pour un fichier de plus de 30~Mio,
        Amcache ne hache que les 31\,457\,280 premiers octets ;
        \champ{Sha1CoversBytes} le signale, et ce SHA-1 ne se compare pas à celui
        du fichier entier.
  \item[\fichier{shimcache.json}] cache de compatibilité : chemin et date de
        modification du fichier, position dans le cache (\champ{CachePosition},
        0 pour la plus récente). Une entrée prouve que le binaire était présent,
        pas qu'il a été exécuté. L'indicateur d'exécution que d'autres outils
        publient n'est \emph{pas} relevé : mesuré contre le Prefetch, son « oui »
        est presque toujours juste, mais son « non » est faux pour 9 binaires
        exécutés sur 10 — il ferait conclure qu'un programme n'a jamais tourné.
  \item[\fichier{bams.json}] dernière exécution par programme et par utilisateur
        (Background Activity Moderator). Les ruches étant lues sur le disque,
        une exécution de la ou des dernières minutes peut manquer : Windows
        garde une modification du registre en mémoire avant de l'écrire (de
        56 à 82~s mesurées sur la machine de test, que \texttt{RegFlushKey}
        n'a pas avancées), et \wac{} ne force pas cette écriture, qui
        modifierait le disque examiné.
  \item[\fichier{userassists.json}] programmes lancés depuis l'interface :
        compteur, nombre et durée de premier plan (\champ{FocusTimeMs}),
        dernière exécution.
  \item[\fichier{muicache.json}] noms d'affichage des programmes exécutés.
  \item[\fichier{run.json}] programmes lancés au démarrage ou à l'ouverture de
        session (clés Run et RunOnce), avec la dernière écriture de la clé : dans
        SOFTWARE (toute la machine), dans la ruche DEFAULT (le compte système) et
        dans le NTUSER.DAT de chaque profil (\champ{Hive}). \champ{Type}
        \texttt{REG\_EXPAND\_SZ} : la commande contient des variables
        (\texttt{\%SystemRoot\%}\ldots), données telles qu'écrites. Aussi :
        les clés Run des programmes 32~bits (\texttt{SOFTWARE (32-bit)}) ;
        \texttt{StartupApproved} (\champ{Key} \texttt{StartupApproved\textbackslash Run},
        \texttt{Run32}, \texttt{StartupFolder}), l'état d'une entrée tel que
        l'utilisateur l'a réglé dans le Gestionnaire des tâches ou les Paramètres :
        \champ{Status}, les 4 premiers octets tels que stockés (leur sens n'est pas
        mesuré, donc pas interprété), et \champ{Changed}, l'instant du
        changement ; \texttt{RunNotification}, une entrée de démarrage que Windows
        a remarquée (\champ{Notification}). Un nom sans valeur Run correspondante
        est une entrée retirée depuis.
  \item[\fichier{persistence.json}] persistance du registre au-delà des clés
        Run : IFEO (\texttt{Debugger}, \texttt{GlobalFlag}), SilentProcessExit,
        AppInit\_DLLs et AppCertDlls, Winlogon (\texttt{Shell},
        \texttt{Userinit}, \texttt{Taskman}, \texttt{Notify}), paquets LSA,
        moniteurs d'impression, \texttt{BootExecute}, Active Setup, assistants
        NetSh, commandes d'ouverture des types exécutables (\texttt{exefile},
        \texttt{batfile}, \texttt{cmdfile}\ldots, machine et utilisateur), et
        détournements COM des classes de chaque utilisateur (UsrClass.dat). Pour chaque valeur : mécanisme (\champ{Mechanism}),
        ruche, clé, nom, type, valeur telle qu'écrite, dernière écriture de la
        clé, et avec \opt{binary} les empreintes du fichier visé. Toutes les
        valeurs sont données, celles de Windows comprises : distinguer
        l'ajouté du défaut revient à l'analyse. Les tâches BITS ne sont pas
        encore lues.
  \item[\fichier{office\_trust\_records.json}] documents Office qu'un utilisateur
        a approuvés (« Activer le contenu ») : application, version, document,
        date ; \champ{MacrosEnabled} quand il a activé les macros — la porte de
        la plupart des infections par macro.
  \item[\fichier{typed\_paths.json}, \fichier{word\_wheel\_query.json}] chemins
        tapés dans la barre d'adresse de l'Explorateur et termes qui y ont été
        cherchés, par utilisateur et par position (\champ{MruPosition}, 0 = le plus
        récent, comme dans \fichier{shellbags.json}).
  \item[\fichier{root\_certificates.json}] certificats des magasins racines de la
        machine et de chaque utilisateur : sujet, émetteur, validité ; une
        racine ajoutée fait accepter tout ce qu'elle signe (interception TLS,
        logiciel signé). \champ{ThumbprintMismatch} signale un certificat dont le
        SHA-1 n'est pas le nom de sa clé.
  \item[\fichier{internet\_settings.json}] proxy et script de configuration
        automatique de chaque utilisateur.
  \item[\fichier{firewall.json}] règles du pare-feu (nom, action, sens, ports,
        programme ; tous les champs dans \champ{Fields}) et réglages des profils
        domaine, privé et public.
  \item[\fichier{installed\_programs.json}] programmes installés, lus dans les
        clés \texttt{Uninstall} de la machine (64 et 32~bits) et de chaque
        utilisateur, et produits Windows Installer (\champ{Scope} \texttt{MSI
        product}, avec le paquet mis en cache \champ{LocalPackage}) : toutes leurs valeurs (nom, éditeur,
        version, dossier, commande de désinstallation, \champ{InstallDate} --- le
        jour seul, tel qu'écrit) et la date d'écriture de la clé. Un outil
        d'accès à distance installé, puis désinstallé, y laisse souvent sa trace.
  \item[\fichier{defender\_configuration.json}] configuration de Microsoft
        Defender lue dans le registre : les \textbf{exclusions} (chemins,
        extensions, processus, adresses ; \champ{Source} \texttt{Settings} ou
        \texttt{Policy}), qu'un intrus ajoute pour que ses outils ne soient pas
        analysés, et les valeurs qui règlent ou coupent la protection (temps
        réel, protection contre la falsification, envoi d'échantillons\ldots).
  \item[\fichier{explorer\_recent\_docs.json}] fichiers ouverts par chaque
        utilisateur, tels que l'Explorateur les retient (\texttt{RecentDocs}) :
        par extension (\champ{Extension}) et pour toutes extensions confondues,
        avec leur rang (\champ{MruPosition}, 0 pour le plus récent), l'élément de
        shell du raccourci créé dans Recent, et la date d'écriture de la liste,
        qui est celle de l'ouverture du fichier en position~0. Le nom survit au
        fichier et au raccourci.
  \item[\fichier{mount\_points.json}] volumes et partages réseau qu'a vus
        l'Explorateur de chaque utilisateur (\texttt{MountPoints2}) : volume
        (GUID, à rapprocher de \fichier{mounted\_device.json} pour une clé USB)
        ou partage (\champ{Share}, \texttt{\textbackslash\textbackslash serveur\textbackslash partage}),
        date d'écriture de l'entrée --- la dernière fois que l'Explorateur l'a
        vue ---, ses valeurs, et ses sous-clés avec leur commande
        (\texttt{Shell\textbackslash AutoRun\textbackslash command} : une
        persistance ancienne).
  \item[\fichier{known\_networks.json}] réseaux auxquels la machine s'est
        connectée (\texttt{NetworkList}) : nom, type, catégorie, adresse
        matérielle de la passerelle et suffixe DNS de chaque signature, première
        et dernière connexion. Ces deux dates sont écrites par Windows en heure
        locale : \wac{} les donne en UTC avec le décalage en vigueur à leur
        date (mesuré contre le journal \texttt{NetworkProfile}). Elles disent où
        la machine a été.
  \item[\fichier{program\_usage.json}] ce que chaque utilisateur a lancé ou
        utilisé, source par source (\champ{Source}) : \texttt{FeatureUsage}
        (lancements, bascules, clics sur la barre des tâches, par application,
        \champ{Count}), \texttt{JumplistData} (dernière utilisation,
        \champ{LastUsed}), le magasin de l'Assistant de compatibilité des
        programmes (programmes exécutés), les modes de compatibilité
        (\texttt{Layers}), les programmes diagnostiqués par RADAR, les éléments
        épinglés à la barre des tâches (\texttt{Taskband}), chacun décodé comme
        un raccourci complet, et les programmes choisis pour ouvrir chaque type
        de fichier (\texttt{FileExts} : \texttt{OpenWithList} avec
        \champ{Extension} et \champ{MruPosition}, 0 pour le plus récent ;
        \texttt{UserChoice}, le programme par défaut). \texttt{OpenWithProgids}
        n'est pas relevé : les installeurs le remplissent, l'utilisateur n'y
        est pour rien.
  \item[\fichier{devices.json}] périphériques que la machine a connus
        (\champ{Source}) : USB et SCSI (\texttt{Enum}), avec leurs valeurs et les
        dates d'installation, de première installation, de dernière arrivée et
        de dernier retrait (\champ{Install}, \champ{FirstInstall},
        \champ{LastArrival}, \champ{LastRemoval}, en UTC) ; périphériques
        portables (téléphone, appareil photo) ; noms de volume que la recherche
        Windows a retenus par lettre de lecteur (\texttt{VolumeInfoCache}) ;
        interfaces de disque, de volume et de périphérique portable
        (\texttt{DeviceClasses}), écrites à leur première apparition.
  \item[\fichier{network\_shares.json}] dossiers que la machine partage
        (\texttt{LanmanServer\textbackslash Shares}) : nom, chemin, réglages tels
        que stockés, et descripteur de sécurité décodé --- propriétaire, et pour
        chaque entrée de la liste de contrôle d'accès son type, son masque, le
        SID et son nom. Les partages administratifs (\texttt{C\$},
        \texttt{ADMIN\$}) n'y figurent pas : Windows ne les stocke pas.
  \item[\fichier{network\_configuration.json}] configuration réseau lue dans
        les ruches : nom d'hôte et domaine (\champ{Global}), et pour chaque carte
        sa description et son pilote, une adresse matérielle forcée
        (\texttt{NetworkAddress}), son adresse matérielle courante et
        d'origine, sa dernière configuration IPv4 et IPv6, et son bail DHCP
        (obtention et fin, en UTC). Elle reste lisible quand la machine est
        débranchée ou le bail expiré, ce que le relevé vivant ne montre plus.
  \item[\fichier{event\_log\_configuration.json}] configuration de chaque
        journal d'événements lue dans les ruches : les canaux
        (\texttt{WINEVT\textbackslash Channels} : \champ{Enabled},
        \champ{Type}, \champ{OwningPublisher}, droits d'accès
        \champ{ChannelAccess}, taille et rétention si elles sont fixées) et les
        journaux classiques (\texttt{Services\textbackslash EventLog} :
        fichier, \champ{MaxSize}, \champ{Retention}, accès) et les sessions
        ETW lancées au démarrage (\texttt{WMI\textbackslash Autologger} :
        \champ{Start}, fichier, taille). Un intrus qui coupe un canal, réduit
        un journal, en ouvre l'accès ou arrête une session comme
        \texttt{EventLog-Security} y laisse sa trace ; \champ{KeyLastWrite} date la dernière modification de chacun.
  \item[\fichier{fat\_volumes.json}] volumes FAT12, FAT16, FAT32 et exFAT, clés
        USB et cartes du suspect comprises (le volume de la collecte excepté) :
        chaque entrée, présente ou \textbf{supprimée} (\champ{Deleted}), son
        chemin, son nom 8.3, sa taille, ses attributs et ses dates. Un nom
        supprimé en FAT garde son nom long s'il n'a pas été écrasé ; sinon son
        nom 8.3, première lettre perdue (\champ{NameIncomplete}). Chaque date
        est donnée en UTC et en heure locale dérivée de l'UTC, comme dans tous
        les artefacts : un même instant se lit partout à la même heure. FAT
        stocke l'heure locale calculée par Windows au moment où il l'écrit : une
        date de modification antérieure à la création est une date conservée
        par une copie, écrite au moment de la copie, et convertie avec le
        décalage de ce moment (\champ{ModifiedWrittenAt} \texttt{Copy},
        \champ{ModifiedOffsetApplied}, \champ{ModifiedAsStored}) — sans quoi un
        fichier d'hiver copié en été serait décalé d'une heure. exFAT garde
        son décalage (\champ{ModifiedOffset}\ldots), la date UTC est exacte. La
        collecte ne copie que les métadonnées du volume (secteur de démarrage,
        FAT, répertoires), pas le contenu des fichiers ; avec \opt{binary}, un
        programme de ces volumes cité par un artefact (lancé depuis une clé,
        par exemple) est empreinté et prélevé comme sur NTFS.
  \item[\fichier{srum.json}] usage des ressources (SRUM, \path|sru\SRUDB.dat|),
        une entrée par ligne de chaque table : octets envoyés et reçus par
        application et par interface (une exfiltration s'y lit, avec le
        programme), cycles processeur, octets lus et écrits, connexions ; 30 à
        60 jours. \champ{Table} (GUID du fournisseur), \champ{Provider} (nom
        documenté : \texttt{NetworkUsage}, \texttt{ApplicationResourceUsage}\ldots),
        puis les colonnes par leur nom (lettres et chiffres seuls), et
        \champ{App}, \champ{UserSID}, \champ{UserName} résolus par
        \texttt{SruDbIdMapTable}. La base est copiée vivante : la dernière
        heure est dans ses journaux \path|SRU*.log|, copiés dans la consigne pour
        un rejeu (\texttt{esentutl /r}) sur une copie.
  \item[\fichier{bits\_jobs.json}] travaux BITS (\path|qmgr.db|) : nom,
        description, type, état, propriétaire, \textbf{commande de
        notification} (lancée à la fin ou à l'échec du transfert : une
        persistance), délais, dates, et chaque fichier : URL source,
        destination, fichier temporaire, tailles. Le blob complet est donné en
        hexadécimal : BITS~10 y garde plus (le programme créateur) à des
        endroits non décrits.
  \item[\fichier{wmi\_persistence.json}] abonnements aux événements WMI, lus
        \textbf{hors ligne} dans le dépôt CIM copié en brut
        (\path|wbem\Repository|) : aucune requête WMI, rien exécuté sur la
        machine. Dans \emph{chaque} espace de noms, chaque filtre (requête
        WQL), consommateur (ligne de commande, script VBScript ou JScript,
        journal, courriel, et toute classe qu'un intrus aurait définie),
        liaison et minuterie, avec toutes ses propriétés (\champ{Properties} :
        nom et valeur ; \champ{CreatorSID} en clair). C'est une persistance
        invisible au registre. Les deux dates d'en-tête d'une instance ne sont
        pas émises : ce ne sont pas des dates de création.
  \item[\fichier{startup\_folders.json}] contenu des dossiers de démarrage, de
        tous les utilisateurs (\champ{Scope} \texttt{AllUsers}) et de chaque
        profil (\texttt{User}, avec \champ{SID}) : chaque fichier, ses dates et
        ses empreintes ; un raccourci est décodé (\champ{Shortcut}, comme dans
        \fichier{recentdocs.json}).
  \item[\fichier{powershell\_history.json}] historique PSReadLine de chaque
        profil et de chaque hôte (\champ{Host} : \texttt{ConsoleHost},
        l'ISE, VS~Code) : les commandes dans l'ordre où elles ont été tapées,
        une commande sur plusieurs lignes réunie. Aucune date par commande :
        \champ{SourceModified} date la dernière.
  \item[\fichier{error\_reports.json}] rapports d'erreurs Windows
        (\fichier{Report.wer}) de la machine et de chaque profil : programme,
        chemin, instant (\champ{EventTimeUtc}), signature du problème
        (module fautif, version), toutes les valeurs du rapport dans
        \champ{Values}.
  \item[\fichier{defender\_log.json}] journaux de support de Microsoft Defender
        (\fichier{MPLog-*.log}) : chaque ligne telle qu'écrite, datée quand elle
        commence par une date UTC. Analyses à l'exécution, détections,
        exclusions, mises à jour.
  \item[\fichier{defender\_detections.json}] une entrée par fichier
        DetectionHistory, qui survit à l'effacement du journal d'événements :
        menace, ressource (fichier, clé), compte, processus, date de détection
        initiale, entrée de quarantaine, et les propriétés nommées
        \champ{ThreatTracking} (empreintes, taille, heure de début). Le format
        n'est pas documenté : seuls les champs dont la position a été confirmée
        contre \texttt{MSFT\_MpThreatDetection} sont nommés ; les autres sont
        rendus tels que lus (\champ{UnnamedFields} : position, type, valeur),
        sans signification inventée.
  \item[\fichier{defender\_quarantine.json}] entrées de la quarantaine :
        menace, instant de mise en quarantaine, chemin d'origine du fichier
        retiré, ses dates de création et de modification, ses attributs, sa
        taille, et l'empreinte de son contenu (\champ{ContentSha256}). La
        quarantaine est copiée \textbf{chiffrée}, telle que Defender la garde :
        rien de malveillant n'est écrit en clair, ni sur le support de collecte
        ni sur le poste d'analyse. Le contenu n'est déchiffré qu'en mémoire,
        pour son empreinte.
  \item[\fichier{hosts.json}] entrées du fichier \fichier{hosts} : adresse et
        noms, numéro de ligne ; un Windows intact n'en a aucune.
\end{description}

\subsection{Fichiers et navigation}

\begin{description}[style=nextline,leftmargin=1.2em]
  \item[\fichier{mft.json}] chaque fichier et répertoire de chaque volume NTFS
        fixe (\champ{Volume} : \texttt{C:}, \texttt{G:}\ldots{} ; un numéro
        d'enregistrement n'a de sens que dans son volume),
        présent ou \textbf{supprimé} (\champ{Deleted}) tant que NTFS n'a pas
        réutilisé son enregistrement : chemin reconstitué, taille, attributs, et
        deux jeux de quatre dates en UTC — \champ{CreatedUtc}\ldots{} (attribut
        \path|$STANDARD_INFORMATION|, ce que montre l'Explorateur et que
        réécrit un outil d'antidatage) et \champ{FileNameCreatedUtc}\ldots{}
        (attribut \path|$FILE_NAME|, écrit par NTFS seul). Une date
        \path|$STANDARD_INFORMATION| antérieure à celle de \path|$FILE_NAME|
        trahit un antidatage. \champ{PathIncomplete} marque un chemin dont un
        répertoire parent a été réutilisé depuis : le chemin n'est pas inventé.
        Un fichier à liens physiques (\champ{HardLinks}) est un seul
        enregistrement sous plusieurs noms — \path|System32\crypt32.dll| et son
        original sous \path|WinSxS| : \champ{HardLinkPaths} donne ses autres
        chemins, chacun cherchable.
  \item[\fichier{usn\_journal.json}] journal des modifications de chaque
        volume NTFS fixe (\path|$UsnJrnl:$J|, \champ{Volume} ; les numéros
        \champ{Usn} croissent au sein d'un volume) : chaque création, écriture, renommage et
        \textbf{suppression}, daté, avec ses motifs (\champ{Reasons}, un bit
        inconnu rendu en hexadécimal) et le chemin reconstitué par la
        \path|$MFT|. Il garde la trace de fichiers déposés puis effacés, que la
        \path|$MFT| a déjà oubliés.
        La \path|$MFT|, le journal des modifications et \path|$LogFile| (journal
        de transactions, non converti) sont en outre copiés bruts dans la
        consigne, pour refaire plus tard la chronologie avec un autre outil.
  \item[\fichier{recentdocs.json}] raccourcis de documents récents : cible,
        horodatages du raccourci et de la cible, taille de la cible, volume
        (type, numéro de série, nom), chemin réseau, liste d'éléments du shell,
        et les blocs de données supplémentaires : le bloc de suivi
        (\champ{Tracker} : nom NetBIOS de la machine où était la cible,
        identifiants de suivi, et ce qu'en révèle le GUID de version 1 — la date
        de début du suivi et l'adresse matérielle de la machine), chemins avec
        variables d'environnement, dossier connu, propriétés de la cible.
  \item[\fichier{jumplistAutomaticDestinations.json},
        \fichier{jumplistCustomDestinations.json}] listes de sauts par
        application, avec les raccourcis qu'elles contiennent ; pour les listes
        automatiques, chaque raccourci est accompagné de son entrée DestList
        (dernier accès, nom de la machine, identifiants de suivi avec leur date
        et leur adresse matérielle, épinglage, nombre d'ouvertures
        \champ{InteractionCount}, rang \champ{MruPosition} — 0 pour le plus
        récent —, nom du dossier connu visé).
  \item[\fichier{shellbags.json}] dossiers parcourus dans l'Explorateur,
        arborescence reconstituée ; pour chaque dossier, son rang parmi ses
        voisins (\champ{MruPosition}, 0 pour le plus récent), la date d'écriture
        de la clé qui le porte (\champ{LastWriteTime}) et de sa propre clé
        (\champ{KeyLastWriteTime} : son ajout, ou le parcours d'un sous-dossier),
        et son \champ{NodeSlot}. Les « première » et « dernière interactions »
        d'autres outils ne sont que des lectures de ces deux dates.
  \item[\fichier{mrus.json}, \fichier{mruApps.json}] listes de fichiers et
        d'applications récemment utilisés.
  \item[\fichier{zone\_identifiers.json}] flux \fichier{Zone.Identifier}
        (« marque du web ») que navigateurs et messageries écrivent sur ce qu'ils
        téléchargent : fichier, zone (\champ{ZoneId}, \champ{ZoneName} —
        \texttt{Internet} pour~3), page d'origine (\champ{ReferrerUrl}) et
        adresse du fichier (\champ{HostUrl}), toutes les valeurs dans
        \champ{Values}. Lus dans la copie de la \path|$MFT|, où ce petit flux
        est résident : aucune lecture de plus sur la machine examinée ; un flux
        non résident, rare, est signalé (\champ{ContentNotResident}).
  \item[\fichier{recycle\_bin.json}] fichiers mis à la corbeille, sur chaque
        volume NTFS fixe, d'après leur index \path|$I| : chemin d'origine,
        taille, instant de la suppression (\champ{DeletedUtc}), et l'utilisateur
        (\champ{SID}). Seuls les \path|$I| sont prélevés ; les fichiers eux-mêmes
        (\path|$R|) figurent dans \fichier{mft.json}.
\end{description}

\subsection{Périphériques}

\begin{description}[style=nextline,leftmargin=1.2em]
  \item[\fichier{setupapi.json}] sections de \fichier{setupapi.dev.log} et de
        ses rotations : opération (\champ{Operation}), périphérique ou pilote
        (\champ{Subject}, avec le numéro de série d'une clé USB), début et fin
        convertis depuis l'heure locale de la machine examinée, résultat. La
        première connexion d'un périphérique s'y date, même quand la clé
        \fichier{USBSTOR} a été nettoyée.
  \item[\fichier{Usbstor.json}] périphériques de stockage USB connus du système :
        \champ{Vendor}, \champ{Product}, \champ{Revision} et \champ{Type}, tirés
        de \champ{DeviceId} ; \champ{SerialNumber}, tiré de \champ{InstanceId}
        — absent quand l'identifiant a été fabriqué par Windows (deuxième
        caractère « \& »), le périphérique n'en déclarant aucun ; dates de
        première et de dernière insertion.
  \item[\fichier{mounted\_device.json}] correspondance entre lettres de lecteur et
        périphériques. \champ{Drive} nomme la lettre ou le volume ;
        \champ{Type} dit la forme de la valeur : \texttt{GPT partition}
        (\champ{PartitionGuid}), \texttt{MBR partition}
        (\champ{DiskSignature} et \champ{PartitionOffset}, en octets),
        \texttt{Device path} (\champ{Device}, chemin du périphérique, avec le
        numéro de série d'une clé USB) ou \texttt{Unknown} (\champ{Data}, les
        octets en hexadécimal).
\end{description}

\subsection{Journaux d'événements (\opt{events})}

\fichier{events.json} contient un objet par événement :

\begin{description}[style=nextline,leftmargin=1.2em]
  \item[\champ{EvtSystem*}] la section système : fournisseur, identifiant,
        niveau, tâche, mots-clés, horodatage de création, numéro d'enregistrement,
        processus et fil, canal, machine, utilisateur ;
  \item[\champ{EvtEventData}] les données de l'événement, en liste de paires
        \champ{Name} / \champ{Value} ;
  \item[\champ{EvtEventMessage}] le message lisible, tel que l'Observateur
        d'événements l'afficherait, reconstitué à partir des ressources du
        fournisseur prélevées sur la machine ;
  \item[\champ{EvtSourceLog}] le fichier journal dont provient l'événement. Un
        même canal peut être porté par plusieurs fichiers (journal courant et
        archives) dont les numéros d'enregistrement se recouvrent ;
  \item[\champ{EvtRecovered}] \texttt{true} pour un événement d'un journal
        \emph{vidé} : Windows remet l'en-tête du bloc à zéro sans effacer les
        enregistrements, que WAC relit. L'Observateur d'événements ne les montre
        plus ; ils restent la trace de ce que l'effacement voulait faire
        disparaître, et ne doivent jamais être lus comme des événements vivants.
\end{description}

\begin{note}[Messages incomplets]
Une marque \texttt{\%1}, \texttt{\%2}\ldots{} restée dans un message signifie que
l'événement ne porte pas la donnée attendue : elle est laissée visible pour que la
phrase ne paraisse pas complète. Tous les fournisseurs ne livrent pas de
ressources de messages : leurs événements n'ont alors pas de
\champ{EvtEventMessage}, mais leurs données restent complètes.
\end{note}

% ============================================================ 6
\chapter{Lire les données}\label{chap:lecture}

\section{Horodatages}

Chaque date est émise deux fois : \champ{X} en heure locale et \champ{XUtc} en
UTC. L'heure locale est celle du \textbf{suspect}, calculée avec le fuseau lu
dans sa ruche \fichier{SYSTEM} — pas celle de la machine de collecte. Les dates
portent la fraction de seconde quand la source la fournit, et seulement dans
la précision de cette source : sept chiffres pour un \texttt{FILETIME}, trois
pour une date à la milliseconde, aucun pour une date à la seconde (date FAT,
précise à deux secondes, date lue en texte, temps Unix).

Chaque date porte le décalage en vigueur \textbf{à cette date}, heure d'été
comprise : une date de janvier relevée en été à Paris est à \texttt{+01:00},
une date de juillet à \texttt{+02:00}. \wac{} applique les règles de la machine
examinée année par année (clé \fichier{Dynamic DST} de la ruche
\fichier{SOFTWARE}) : une date de 2005 aux États-Unis suit les règles de 2005.
L'heure locale écrite est toujours dérivée de l'heure UTC : un même instant se
lit à la même heure dans tous les artefacts. Deux dates n'ont aucun fuseau
connaissable et sont données \textbf{telles qu'écrites}, sans décalage ni UTC,
avec \champ{TimeZone} qui le dit : la date d'un élément vu dans une archive
(\champ{ModifiedAsWritten} ; une entrée ZIP garde l'heure d'horloge de la
machine qui a fait l'archive) et une date OLE d'un magasin de propriétés
(\champ{AsWritten}). Les placer sur une chronologie demande de connaître
leur origine : \wac{} ne l'invente pas. Pour une valeur que l'artefact
enregistre en heure locale (volume FAT, \fichier{setupapi.dev.log}), l'heure répétée à l'automne est lue comme sa première occurrence, et
l'heure sautée au printemps avec le décalage d'avant le changement, comme le
fait Windows.

\section{Noms des comptes}

Les champs qui nomment un compte (\champ{SIDName}, \champ{Owner},
\champ{RunAs}) sont tirés des preuves seules, sans interroger le système ni le
réseau : comptes et groupes locaux de la base SAM (dans la langue de la
machine examinée), SID bien connus sous leur nom canonique
(\texttt{SYSTEM}, \texttt{LOCAL SERVICE}\ldots), SID de services
(\texttt{NT SERVICE\textbackslash\ldots}), et dossier de profil pour un
compte de domaine. Un SID qu'aucune de ces sources ne nomme est publié sans
nom : le SID reste la preuve, à rapprocher du domaine par l'analyste.

\section{Noms des GUID}

Les champs qui nomment un GUID (classe COM, dossier connu, gestionnaire d'une
tâche) viennent d'abord d'une table de référence, puis des ruches de la machine
examinée : ses classes (\fichier{SOFTWARE\textbackslash Classes\textbackslash
CLSID}), ses dossiers connus (\fichier{FolderDescriptions}), et les classes de
chaque utilisateur (\fichier{UsrClass.dat}). Un composant installé sur la
machine y est donc nommé, avec le module qui le sert. Un GUID que rien ne nomme
est publié sans nom.

\section{Un champ absent n'est pas un champ vide}

Une valeur sans source n'est pas émise du tout. Une clé vide se lirait comme une
lecture ratée, et ne se distinguerait pas de « la ruche ne contient pas cette
valeur ». Ce qui n'a pas pu être lu est consigné dans
\fichier{investigation.json}.

\section{Artefact non collecté}

Quand un artefact entier n'a pas pu être lu, son fichier le dit explicitement :

\begin{lstlisting}
{
  "Artifact": "USBSTOR (registre)",
  "CollectionStatus": "NotCollected",
  "Error": "0x02 Le fichier spécifié est introuvable.",
  "Note": "La lecture de cet artefact a échoué : l'absence de données
           ci-dessus ne signifie PAS qu'aucune trace n'existe sur le système."
}
\end{lstlisting}

\begin{attention}
Un fichier d'artefact vide et un fichier « NotCollected » ne disent pas la même
chose. Le premier indique qu'aucune entrée n'a été trouvée ; le second, que
la source n'a pas pu être lue.
\end{attention}

\section{Données non décodées}

Un élément que \wac{} ne sait pas décoder (élément du shell inconnu, bloc
d'extension non reconnu) est accompagné de ses octets bruts en hexadécimal, pour
qu'un analyste puisse le décoder plus tard. Avec \opt{dump}, la valeur brute des
shellbags et des raccourcis est ajoutée systématiquement.

\section{Chemins}

Les chemins et valeurs sont restitués \textbf{tels qu'ils figurent sur le
disque}, sans transformation : l'échappement n'intervient qu'à l'écriture du
JSON.

% ============================================================ 7
\chapter{Traces laissées sur la machine examinée}\label{chap:traces}
\markboth{Traces sur la machine examinée}{}

À lire avant d'employer \wac{} sur une affaire. Ce chapitre énumère ce que \wac{}
fait à la machine, opération par opération, y compris ce qu'il ne peut pas
éviter.

\section{Ce que WAC fait encore}

\begin{longtable}{@{}L{4.3cm}L{10.1cm}@{}}
\toprule
\bfseries Opération & \bfseries Trace \\
\midrule
\endhead
Lecture directe du volume & Un accès par volume effectivement lu (en général un
  seul). Aucun fichier ouvert, donc aucune date d'accès modifiée. Si l'audit
  d'accès aux objets est activé, l'ouverture du volume peut être journalisée
  (Security 4656/4663). \\
Écriture des résultats & Sur le support de collecte uniquement. Rien n'est écrit
  sur le disque examiné. \\
Rejeu des journaux de ruche & Sur la copie de travail uniquement, avec journal
  d'annulation. \\
État des services & Une requête en lecture au gestionnaire de services. Aucun
  service n'est ouvert ni modifié. \\
Processus & Un instantané du noyau ; aucun processus n'est ouvert. Relevé des
  propriétaires : une requête au service Terminal Services. \\
Sessions & Une requête en lecture à LSASS. \\
État réseau & Lectures en mémoire des tables de la pile TCP/IP : aucun paquet
  émis, aucun nom résolu. \texttt{iphlpapi.dll} n'est chargée que si
  \texttt{artefacts.network} est activé. \\
Heure de démarrage & Une requête en mémoire au noyau. \\
Journaux d'événements (\opt{events}) & Lus par le même accès direct au volume :
  le service de journalisation n'est pas sollicité. \\
Fichiers cités (\opt{binary}) & Lus par le même accès direct. Aucun fichier
  ouvert. \\
Vérification d'authenticité (\opt{binary}) & Catalogues lus par le même accès
  direct ; vérification en mémoire. Aucun service sollicité, aucun accès au
  magasin de certificats, aucun accès réseau. \\
\bottomrule
\end{longtable}

\section{Ce que WAC ne fait plus}

Les versions antérieures créaient un cliché instantané (VSS), passaient par COM et
WMI, interrogeaient le planificateur de tâches et le service de journalisation,
ouvraient chaque processus et chaque service, lisaient la liste des profils dans
le registre vivant et ouvraient chaque fichier cité pour en calculer
l'empreinte. Toutes ces opérations, qui laissaient des entrées dans les journaux
ou modifiaient des dates d'accès, ont été supprimées.

\section{Ce qu'aucun outil ne peut éviter}

Ces traces sont listées pour que l'analyste reconnaisse l'empreinte de \wac{} au
lieu de la prendre pour l'activité étudiée :

\begin{itemize}
  \item \textbf{lancer \fichier{WAC.exe} est une exécution} : elle produit un
        fichier Prefetch pour \fichier{WAC.exe}, une entrée Amcache et, si
        l'audit de création de processus est actif, un événement Security 4688.
        Attendez-vous à retrouver \wac{} dans les artefacts qu'il collecte ;
  \item \textbf{brancher la clé USB} l'inscrit dans \fichier{USBSTOR},
        \fichier{MountedDevices} et les journaux associés ;
  \item \textbf{ouvrir une session} pour lancer l'outil crée une session et ses
        entrées de journal ;
  \item \wac{} \textbf{occupe de la mémoire} pendant la collecte, ce qui peut
        évincer des pages vers le fichier d'échange. Sa lecture du volume, elle,
        contourne le cache (lecture non tamponnée) : les gigaoctets lus ne
        remplacent pas le cache de fichiers de la machine.
\end{itemize}

% ============================================================ 8
\chapter{Limites connues}

\begin{description}[style=nextline,leftmargin=1.2em]
  \item[Volumes ni NTFS ni FAT] Un fichier situé sur un CD-ROM n'est pas lu :
        il est signalé comme non lu.
  \item[Variables propres à un utilisateur] Un chemin exprimé avec
        \path|%APPDATA%|, \path|%LOCALAPPDATA%| ou \path|%USERPROFILE%| (dans une
        tâche planifiée, par exemple) n'est pas résolu : le compte concerné n'est
        pas connu à cet endroit, et deviner donnerait l'empreinte d'un autre
        fichier.
  \item[Exécutable supprimé] Un processus dont l'exécutable a été effacé du disque
        n'a pas d'empreinte : il ne reste que la mémoire, hors du périmètre de
        \wac{}.
  \item[Messages d'événements] Environ un tiers des événements reçoivent un
        message lisible sur la machine de test ; les autres fournisseurs ne livrent
        pas de ressources exploitables. Les données, elles, sont toujours
        complètes.
  \item[Signés mais non authentifiés] Quelques sortes de fichiers réellement
        signés restent prélevées, dans le sens prudent : les applications du
        Microsoft Store (signées au niveau du paquet, et non fichier par fichier)
        et les scripts Windows Script Host signés seulement dans leur texte.
\end{description}

% ============================================================ 9
\chapter{Dépannage}

\begin{longtable}{@{}L{5cm}L{9.4cm}@{}}
\toprule
\bfseries Message ou constat & \bfseries Cause et remède \\
\midrule
\endhead
Échec de « Check administrator rights » & L'invite n'a pas été lancée en tant
  qu'administrateur. Relancez-la avec les droits d'administrateur. \\
Échec de « Check OS >= Windows 10 » & Système antérieur à Windows~10 : non pris
  en charge. \\
Échec de « Checking the collection medium », répertoire non vide & Le dossier de
  sortie contient une collecte antérieure. Indiquez un dossier neuf avec
  \opt{output}. \\
Échec de « Checking the collection medium », place insuffisante & Le support
  est trop petit pour les options choisies (section~\ref{sec:place}). Utilisez un
  support plus grand ou retirez \opt{binary}. \\
Un fichier d'artefact porte \texttt{"CollectionStatus": "NotCollected"} &
  La source n'a pas pu être lue ; le motif est dans \champ{Error}. Les autres
  artefacts ne sont pas affectés. \\
Des pièces en échec au manifeste & Fichiers cités absents du disque (logiciel
  désinstallé, profil de service sans ruche), volume non NTFS\ldots{} Le motif est
  dans \champ{Result}. \\
Le journal indique des binaires « hachés sans être copiés » & La place a manqué
  pendant le prélèvement : les empreintes sont complètes, les copies non. \\
Besoin de comprendre un échec de lecture & Relancez avec \opt{loglevel=2} (et, pour
  la lecture NTFS, \opt{debug}) et consultez \fichier{WAC.exe.log}. \\
\bottomrule
\end{longtable}

% ============================================================ APPENDICES
\appendix

\chapter{Glossaire}

\begin{description}[style=nextline,leftmargin=1.2em]
  \item[Artefact] Trace laissée par l'activité du système ou d'un utilisateur, et
        exploitable en investigation.
  \item[Consigne] Répertoire des copies brutes, scellé, jamais réécrit.
  \item[Empreinte] Valeur calculée sur le contenu d'un fichier (MD5, SHA-1,
        SHA-256) : deux contenus différents ont des empreintes différentes.
  \item[Journal de transaction] Fichier \fichier{.LOG1}/\fichier{.LOG2} contenant
        les modifications récentes d'une ruche, pas encore écrites dans celle-ci.
  \item[Manifeste] Fichier qui identifie chaque pièce de la consigne.
  \item[\path|$MFT|] Table des fichiers d'un volume NTFS.
  \item[Ruche] Fichier constituant une partie de la base de registre
        (\fichier{SYSTEM}, \fichier{SOFTWARE}, \fichier{ntuser.dat}\ldots).
  \item[Sceau] Empreinte du manifeste, conservée à part : elle révèle toute
        retouche du manifeste.
  \item[Travail] Répertoire des copies de travail, sur lesquelles portent l'analyse
        et les éventuelles modifications.
  \item[WOF / Compact OS] Mode de stockage compressé des fichiers système de
        Windows~10 et~11.
\end{description}

\end{document}
